% !TEX program = xelatex
% 数学物理方法2 · §8.4 球坐标下的分离变量（精讲讲义，2026-11-17）
% 依据课件《Chapter 8 - 20261117.pdf》10 页（p1--p10）
\documentclass[11pt]{ctexart}
%=====================================================================
% theme-violet.tex —— 讲义精讲 · 紫花主题
% 用法：\documentclass[11pt]{ctexart}
%       \input{assets/theme-violet}
%       \begin{document} ... \end{document}
% 只用 ctexart 自带能力 + xcolor/enumitem/fancyhdr/pifont/amsmath，
% 不依赖 tikz/tcolorbox，保证 MiKTeX 可直接编译。
%=====================================================================

\usepackage[a4paper,left=2.0cm,right=2.0cm,top=2.1cm,bottom=2.3cm,%
            headsep=12pt,footskip=22pt]{geometry}
\usepackage{amsmath,amssymb}
\usepackage{needspace}
\usepackage{xcolor}
\usepackage{array}
\usepackage{enumitem}
\usepackage{pifont}
\usepackage{fancyhdr}
\usepackage[hidelinks]{hyperref}

%---------------------------------------------------------------- 配色
\definecolor{violet}{RGB}{107,78,143}      % 主色：紫
\definecolor{violetdeep}{RGB}{63,44,92}    % 深紫：标题/强调
\definecolor{petal}{RGB}{183,159,208}      % 花瓣：装饰线、花饰
\definecolor{bandviolet}{RGB}{247,243,251} % 浅紫底：档案/提示框
\definecolor{bandcream}{RGB}{252,249,242}  % 米白底：易错/注意框
\definecolor{rulepurple}{RGB}{217,207,230} % 分隔线
\definecolor{inkgray}{RGB}{92,96,110}      % 次要文字

%---------------------------------------------------------------- 字体
% 有 Noto CJK 用 Noto，没有就退回 ctex 默认（Windows 下为宋体/黑体）
\IfFontExistsTF{Noto Serif CJK SC}{\setCJKmainfont{Noto Serif CJK SC}}{}
\IfFontExistsTF{Noto Sans CJK SC}{\setCJKsansfont{Noto Sans CJK SC}}{}

%---------------------------------------------------------------- 版式
\setlist{nosep,leftmargin=1.9em}
\linespread{1.28}
\setlength{\parskip}{3.5pt}
\setlength{\parindent}{0pt}
\setlength{\tabcolsep}{5pt}
\renewcommand{\arraystretch}{1.28}

%---------------------------------------------------------------- 花饰
% \ding{95} = 五瓣小花，\ding{96} = 星形花；都来自 pifont，必定存在
\newcommand{\orn}{\textcolor{petal}{\ding{95}}}          % 行内小标
\newcommand{\ornstar}{\textcolor{petal}{\ding{96}}}      % 星形花
\newcommand{\flowerrule}{%                                % 居中花饰分隔线
  \par\vspace{7pt}%
  \noindent\makebox[\linewidth]{\color{petal}%
    \rule{0.20\linewidth}{0.6pt}\hspace{8pt}{\small\ding{95}}\hspace{4pt}%
    \raisebox{1.1pt}{\tiny\ding{95}}\hspace{4pt}{\small\ding{95}}\hspace{8pt}%
    \rule{0.20\linewidth}{0.6pt}}%
  \par\vspace{7pt}}

%---------------------------------------------------------------- 标题与节头
\newcommand{\doctitle}[3]{%   #1 主标题  #2 副标题  #3 元信息（可用 \\ 分行）
  \begin{center}
    {\color{petal}\small \orn\enspace #2\enspace \ornstar}\\[7pt]
    {\LARGE\bfseries\color{violetdeep} #1}\\[10pt]
    {\color{petal}\rule{0.55\linewidth}{1pt}}\\[4pt]
    {\small\color{inkgray} #3}
  \end{center}\vspace{2pt}}

\newcommand{\hdr}[2]{%   #1 章节号（如 §0）  #2 标题
  \par\vspace{12pt}%
  \noindent\colorbox{violetdeep}{\color{white}\bfseries\small\;#1\;}\hspace{7pt}%
  {\color{violetdeep}\large\bfseries #2}\par\nobreak
  \vspace{3pt}{\color{petal}\hrule height 1.2pt}\vspace{6pt}}

\newcommand{\subhdr}[1]{%
  \par\vspace{6pt}%
  \noindent{\color{violet}\bfseries #1}\par\vspace{1pt}}

%---------------------------------------------------------------- 提示框
\newcommand{\callout}[1]{% 浅紫底：讲义档案、读法说明
  \par\medskip\noindent
  \fcolorbox{rulepurple}{bandviolet}{\parbox{\dimexpr\linewidth-2\fboxsep-2\fboxrule}{\vspace{2pt}#1\vspace{2pt}}}%
  \par\medskip}

\newcommand{\warnbox}[1]{% 米白底：易错、边界条件
  \par\medskip\noindent
  \fcolorbox{rulepurple}{bandcream}{\parbox{\dimexpr\linewidth-2\fboxsep-2\fboxrule}{\vspace{2pt}#1\vspace{2pt}}}%
  \par\medskip}

\newcommand{\gapitem}[2]{% #1 缺口类型（紫） #2 内容
  \par\vspace{2pt}\noindent
  \hangindent=1.4em\hangafter=1
  {\color{petal}\small\ding{96}}\enspace{\bfseries\color{violetdeep}#1}\enspace #2\par}

% 题干块（\qstem）：每道例题／作业的「详解」开头必须先复述题干与出处。
% 起因：2026-10-25 第五章讲义把七道作业的题面集中放在一页、详解另起一页，
% 读者翻到「第 4 题」时看不到条件，等于一道没有题干的题目。
% 规则：详解块内自带题干，读者不翻页就能读懂题在问什么。
\newcommand{\qstem}[1]{%
  \par\medskip\noindent
  \fcolorbox{violet}{bandviolet}{%
    \parbox{\dimexpr\linewidth-2\fboxsep-2\fboxrule}{%
      \vspace{2pt}%
      {\small\bfseries\color{violet}题干}\hspace{0.8em}#1%
      \vspace{2pt}}}%
  \par\medskip}

% v2: 保护题干与开头，长题可跨页，续页补上下文。
\newcommand{\qbreak}{\needspace{6\baselineskip}}
% 独立练习：ID、来源、题干、考点、答案标签（通常 ans:ID）。
\newcommand{\practice}[5]{%
  \qbreak\subhdr{练习 #1}\qstem{（#2）#3}%
  \textbf{考点}：#4\par
  详解见附录第~\pageref{#5}~页。\par\medskip}
% 附录开头：ID、来源、条件摘要；自动产生 ans:ID 标签。
\newcommand{\answerstart}[3]{%
  \qbreak\subhdr{#1\quad 详解}\label{ans:#1}\qstem{（#2）#3}}
\newcommand{\answercontinue}[2]{%
  \needspace{5\baselineskip}\subhdr{#1\quad 续解}%
  {\small\color{inkgray}本页接续：#2}\par\medskip}
% 延迟提供默认值，让旧讲义可在导言区保留自己的 \newcommand{\step}。
\AtBeginDocument{\providecommand{\step}[1]{\needspace{3\baselineskip}\par\medskip\textbf{#1}\par}}

% 前提旁注：只放在「需要理解的那一步」后面，给陈述 + 本题代入，不重证。
\newcommand{\aside}[1]{%
  \par\vspace{1pt}\noindent
  \hangindent=1.6em\hangafter=1
  {\small\color{inkgray}%
    \hspace*{0.6em}{\color{petal}$\triangleright$}\hspace{2pt}#1}%
  \par\vspace{2pt}}

%---------------------------------------------------------------- 表格
% 本机 MiKTeX 的 colortbl 与 array 版本不匹配（\rowcolor 会报错），
% 故表格统一用「上粗紫线 + 表头 + 中细线」的三线式，不依赖 colortbl。
\newcommand{\tabtop}{\par\vspace{3pt}{\color{violet}\hrule height 1.2pt}\vspace{1pt}}
\newcommand{\tabmid}{\par\vspace{1pt}{\color{petal}\hrule height 0.5pt}\vspace{1pt}}
\newcommand{\tabbot}{\par\vspace{3pt}{\color{rulepurple}\hrule height 0.5pt}\vspace{8pt}}
\newcommand{\thead}[1]{{\bfseries\color{violetdeep}#1}}

%---------------------------------------------------------------- 页眉页脚
\pagestyle{fancy}
\fancyhf{}
\renewcommand{\headrulewidth}{0.4pt}
\renewcommand{\headrule}{\hbox to\headwidth{\color{rulepurple}\leaders\hrule height \headrulewidth\hfill}}
\fancyhead[L]{\small\color{inkgray}\songti \doctag}
\fancyhead[R]{\small\color{inkgray}\songti \coursetag}
\fancyfoot[C]{\small\color{inkgray}\orn\enspace{\color{violet}\thepage}\enspace\orn}

% 每篇讲义在导言区设这两个值
\newcommand{\doctag}{}
\newcommand{\coursetag}{}
\newcommand{\settags}[2]{\renewcommand{\doctag}{#1}\renewcommand{\coursetag}{#2}}

%---------------------------------------------------------------- 数学记号
\newcommand{\vect}[1]{\boldsymbol{#1}}
\newcommand{\dd}{\,\mathrm{d}}
\newcommand{\ddx}{\frac{\mathrm{d}}{\mathrm{d}x}}
\addtolength{\jot}{2pt}



\usepackage{tikz}
\usetikzlibrary{arrows.meta,calc,positioning,decorations.pathmorphing}

\settags{\S8.4 · 球坐标下的分离变量}{数学物理方法2}

% 同一张表内的中隔线
\newcommand{\tabmidrule}{\noalign{\vspace{1pt}{\color{petal}\hrule height 0.5pt}\vspace{1pt}}}

% 圈码交给中文字体
\xeCJKDeclareCharClass{CJK}{"2460 -> "24FF}

% 表头不进第二行
\renewcommand{\thead}[1]{\leavevmode{\bfseries\color{violetdeep}#1}}
% 文本列：左对齐，避免窄列被拉出大片空隙
\newcolumntype{L}[1]{>{\raggedright\arraybackslash}p{#1}}

% 图与图注不分页
\newenvironment{figbox}{\par\begin{center}\begin{minipage}{0.96\linewidth}\centering}{\end{minipage}\end{center}\par}
\newcommand{\figcap}[1]{\par\vspace{4pt}{\small\color{inkgray}#1}\par}

% 详解块起页保护
\renewcommand{\qbreak}{\needspace{0.28\textheight}}

% 三级标题
\newcommand{\hdrsub}[2]{%
  \par\vspace{10pt}\needspace{4\baselineskip}%
  \noindent
  \colorbox{petal}{\color{white}\bfseries\small\;#1\;}\hspace{7pt}%
  {\color{violetdeep}\bfseries\large #2}\par\nobreak
  \vspace{2pt}{\color{rulepurple}\hrule height 0.6pt}\vspace{5pt}}

% 标题不落页底
\let\origsubhdr\subhdr
\renewcommand{\subhdr}[1]{\par\needspace{3\baselineskip}\origsubhdr{#1}}
\let\orighdr\hdr
\renewcommand{\hdr}[2]{\par\needspace{6\baselineskip}\orighdr{#1}{#2}}

% 定理与证明
\newcommand{\thm}[2]{\par\medskip\noindent{\color{violetdeep}\bfseries #1}\hspace{0.5em}{\bfseries #2}\par\nobreak\vspace{2pt}}
\newcommand{\pfline}{\par\smallskip\noindent\textbf{证明。}\enspace}
\newcommand{\qedsym}{\hfill$\blacksquare$}

% 单位
\newcommand{\unit}[1]{\ifmmode\mathrm{#1}\else$\mathrm{#1}$\fi}
% 组标题
\newcommand{\grp}[1]{%
  \par\vspace{11pt}\needspace{3\baselineskip}%
  \noindent{\color{violetdeep}\bfseries\large\orn\enspace #1}\par\nobreak
  \vspace{3pt}{\color{petal}\hrule height 1pt}\vspace{5pt}}

% 公式间距
\setlength{\headheight}{22pt}
\setlength{\emergencystretch}{1.5em}
\setlength{\abovedisplayskip}{7pt plus 2pt minus 2pt}
\setlength{\belowdisplayskip}{7pt plus 2pt minus 2pt}
\setlength{\abovedisplayshortskip}{4pt plus 1pt}
\setlength{\belowdisplayshortskip}{5pt plus 1pt minus 1pt}

\begin{document}

\doctitle{\S8.4\quad 球坐标下的分离变量}
{从球坐标的 Laplace 算子，到 Legendre 多项式、球谐函数与球贝塞尔}
{依据课件《Chapter 8 -- 20261117.pdf》10 页（\S8.4，p1--p10）\\2026-11-17\quad 数学物理方法2}

\callout{\textbf{这本讲义怎么用。}正文按课件的 10 页推进，凡是课件压缩掉的中间步都补出来，并在小标题里标了页码。三件事请先知道：

\textbf{一、}课件 p1--p2 用两页做完「径向与角向分离、角向再分离、周期条件定 $m$」三件事，跳步很多；本讲把每一步的「为什么」写全（\S1--\S3），因为后面 $\theta$ 方程与径向方程的全部常数都从这三步来。

\textbf{二、}课件 p3--p4 把 Legendre 方程与关联 Legendre 方程直接当已知（只标「Chapter 12」）。但「$\lambda$ 为什么只能是 $l(l+1)$」「$P_l^m$ 里的 $m$ 为什么只能取到 $l$」这两句，不知道有界性的论证就只是口号。本讲补了 \S4：Legendre 方程的有界性、前几个多项式、Rodrigues 公式、正交性与模。

\textbf{三、}课件 p8 的球贝塞尔同样只给结论（标「Chapter 13」）。本讲补了 \S7.3：$x=kr$ 与 $R=x^{-1/2}y$ 两步换元怎样把球贝塞尔方程变成半整数阶贝塞尔方程，以及 $j_0,n_0,j_1,n_1$ 的显式——球对称问题（$l=0$）只用 $j_0,n_0$ 就能手算。}

\clearpage

\subhdr{全书主线}

\begin{center}
\begin{minipage}{0.99\linewidth}
\centering
\begin{tikzpicture}[x=1cm,y=1cm,
  core/.style={draw=violetdeep,fill=bandviolet,rounded corners=3pt,align=left,inner xsep=6pt,inner ysep=5pt,line width=1pt,font=\footnotesize},
  box/.style={draw=violet,fill=bandviolet,rounded corners=2.5pt,align=left,inner xsep=5pt,inner ysep=4pt,line width=0.7pt,font=\footnotesize},
  cream/.style={draw=violet,fill=bandcream,rounded corners=2.5pt,align=left,inner xsep=5pt,inner ysep=4pt,line width=0.7pt,font=\footnotesize},
  ova/.style={-{Stealth[length=2.2mm,width=1.6mm]},violet,line width=0.9pt}]

\node[core,text width=15.2cm,anchor=north west] (CORE) at (0.2,0) {%
\textbf{\color{violetdeep}要解决的事}\quad 球域（$r\le r_0$）、球外或球壳上的 Laplace、波动、输运方程；边界落在 $r=$ 常数与 $\theta=0,\pi$ 上。\\
\textbf{难点}\quad 角向有两个变量 $(\theta,\varphi)$，且 $\nabla^2$ 里的 $\sin\theta$、$\sin^2\theta$ 让 $\theta$ 方程变系数——这一步没有初等解。};

\node[box,text width=7.3cm,minimum height=3.4cm,anchor=north west] (A) at (0.2,-2.10) {%
\textbf{\color{violetdeep}\S1--\S2\quad 第一次分离：$u=R(r)Y(\theta,\varphi)$}\\[1pt]
$\dfrac1R\dfrac{d}{dr}\Big(r^2\dfrac{dR}{dr}\Big)=\lambda$（径向）\\
角向：$\dfrac{1}{Y\sin\theta}\partial_\theta(\sin\theta\,\partial_\theta Y)+\dfrac{1}{Y\sin^2\theta}\partial_\varphi^2Y=-\lambda$\\[1pt]
{\color{inkgray}两边一个是 $r$、一个是 $(\theta,\varphi)$，只能同等于常数}};

\node[box,text width=7.3cm,minimum height=3.4cm,anchor=north west] (B) at (7.9,-2.10) {%
\textbf{\color{violetdeep}\S3\quad 第二次分离：$Y=\Theta(\theta)\Phi(\varphi)$}\\[1pt]
$\Phi''+\mu\Phi=0$，周期条件 $\Rightarrow\mu=m^2$\\
$\sin\theta\dfrac{d}{d\theta}\Big(\sin\theta\dfrac{d\Theta}{d\theta}\Big)+(\lambda\sin^2\theta-m^2)\Theta=0$\\[1pt]
{\color{inkgray}$\mu=m^2$ 是周期条件逼出来的，不是猜的}};

\node[cream,text width=7.3cm,minimum height=3.4cm,anchor=north west] (C) at (0.2,-5.75) {%
\textbf{\color{violetdeep}\S4\quad $x=\cos\theta$：关联 Legendre 方程}\\[1pt]
$\dfrac{d}{dx}\Big[(1-x^2)y'\Big]+\Big(\lambda-\dfrac{m^2}{1-x^2}\Big)y=0$\\
$x=\pm1$ 有界 $\Rightarrow\lambda=l(l+1)$，$l=0,1,2,\dots$\\
解 $P_l^m(\cos\theta)$，$0\le m\le l$\\[1pt]
{\color{inkgray}补：有界性论证、$P_l$ 前几个、正交模方}};

\node[cream,text width=7.3cm,minimum height=3.4cm,anchor=north west] (D) at (7.9,-5.75) {%
\textbf{\color{violetdeep}\S5\quad 径向：Euler 方程}\\[1pt]
$r^2R''+2rR'-l(l+1)R=0$\\
$\Rightarrow R=C_l r^l+D_l r^{-(l+1)}$\\
球内（含 $r=0$）丢 $r^{-(l+1)}$\\[1pt]
{\color{inkgray}三种区域各丢哪一项，只看 $r=0$ 与 $r\to\infty$}};

\node[box,text width=15.2cm,minimum height=3.2cm,anchor=north west] (E) at (0.2,-9.45) {%
\textbf{\color{violetdeep}\S6\quad 球坐标 Laplace 的通解与定系数}\quad
$u=\sum_{l=0}^{\infty}\sum_{m=0}^{l}\big(A_l^m r^l+C_l^m r^{-(l+1)}\big)P_l^m(\cos\theta)\big(\text{$\cos m\varphi$ 或 $\sin m\varphi$}\big)$\\
球面 $u(r_0,\theta,\varphi)=f(\theta,\varphi)$：用 $P_l^m$ 的正交性逐阶读出系数；轴对称（$m=0$）只剩 $P_l(\cos\theta)$\\
{\color{inkgray}教学小例：球内 $u(r_0,\theta)=U_0\cos\theta$ 的解是 $\dfrac{U_0}{r_0}\,r\cos\theta=\dfrac{U_0}{r_0}z$——一个均匀场}};

\node[box,text width=15.2cm,minimum height=2.9cm,anchor=north west] (F) at (0.2,-13.05) {%
\textbf{\color{violetdeep}\S7\quad 波动与输运：球贝塞尔与三重级数}\quad
$u=T(t)V(\mathbf r)\Rightarrow\nabla^2V+k^2V=0$，$T''+a^2k^2T=0$\\
径向 $\big(r^2R'\big)'+\big[r^2k^2-l(l+1)\big]R=0$，换元 $x=kr,\ R=x^{-1/2}y$ $\Rightarrow$ 半整数阶贝塞尔，解 $R=Cj_l(kr)+Dn_l(kr)$\\
球面齐次 $\Rightarrow k=k_{ln}$；三重级数 $\sum_n\sum_l\sum_m R_lP_l^m\Phi_mT_n$};

\node[cream,text width=15.2cm,minimum height=1.4cm,anchor=north west] (G) at (0.2,-16.35) {%
\textbf{\color{violetdeep}\S8\quad 本征值问题总表（课件 p10 的作业）}\quad
直角、极、柱、球四种坐标 $\times$ Laplace、波动、输运三类方程，本征值分别由哪一组条件定：周期、自然（$r=0$ 或 $x=\pm1$ 有界）、齐次边界。};

\draw[ova] ([xshift=-3.4cm]CORE.south) -- ([xshift=-3.4cm]A.north);
\draw[ova] ([xshift=3.4cm]CORE.south) -- ([xshift=3.4cm]B.north);
\draw[ova] ([xshift=-3.4cm]A.south) -- ([xshift=-3.4cm]C.north);
\draw[ova] ([xshift=3.4cm]B.south) -- ([xshift=3.4cm]D.north);
\draw[ova] ([xshift=-3.4cm]C.south) -- ([xshift=-3.4cm]E.north);
\draw[ova] ([xshift=3.4cm]D.south) -- ([xshift=3.4cm]E.north);
\draw[ova] (E.south) -- (F.north);
\draw[ova] (F.south) -- (G.north);

\end{tikzpicture}
\figcap{图 0-1\quad 本节主线：两次分离定出 $\lambda$ 与 $m$，$x=\cos\theta$ 把 $\theta$ 方程变成关联 Legendre 方程，自然边界定 $l$，径向按区域取舍，最后用球谐正交性定系数。}
\end{minipage}
\end{center}

\clearpage

\subhdr{记号}

\begin{center}
\tabtop
\begin{tabular}{@{}L{0.16\linewidth}L{0.76\linewidth}@{}}
\thead{记号} & \thead{含义} \\
\tabmidrule
$r,\ \theta,\ \varphi$ & 球坐标：$0\le r<\infty$，$0\le\theta\le\pi$（自 $+z$ 轴起算的极角），$0\le\varphi<2\pi$（方位角） \\
$x$ & $x=\cos\theta\in[-1,1]$；$\theta=0,1$ 对应 $x=1,-1$ \\
$r_0$ & 球面半径；球内的区域是 $r<r_0$ \\
$\lambda$ & 第一次分离的常数（径向/角向）；后文证明它只能取 $l(l+1)$ \\
$\mu,\ m$ & 第二次分离的常数；周期条件给出 $\mu=m^2$，$m=0,1,2,\dots$ \\
$l,\ m$ & 角量子数与磁量子数；$l=0,1,2,\dots$，$m=0,1,\dots,l$ \\
$P_l,\ P_l^m$ & Legendre 多项式、关联 Legendre 函数（本讲不写 Condon--Shortley 的 $(-1)^m$） \\
$Y_l^m$ & 球谐函数；实形式取 $P_l^m(\cos\theta)\cos m\varphi$ 与 $P_l^m(\cos\theta)\sin m\varphi$ \\
$k$ & 时间分离的常数（波数，单位 $1/$长度）；$k=k_{ln}$ 由球面条件定 \\
$j_l,\ n_l$ & 第一类、第二类球贝塞尔函数；$j_0=\sin x/x$，$n_0=-\cos x/x$ \\
$a$ & 波速或 $\sqrt{\kappa/(\rho c)}$（输运），量纲随方程而定 \\
$A_l^m,\ C_l^m$ & $r^l$ 与 $r^{-(l+1)}$ 的系数；$B_l^m,\ D_l^m$ 是 $\sin m\varphi$ 那一支 \\
\end{tabular}
\tabbot
\end{center}

\warnbox{\textbf{两个记号坑。}①本讲的 $m$ 是整数（$0\le m\le l$），不是 §8.3 柱坐标里那个「$\mu$ 的三种符号」的 $\mu$；两处都写作 $m$ 时看上下文。②$P_l^m$ 的符号约定各家不同：有的教材带 $(-1)^m$（Condon--Shortley），本讲按 $P_l^m=(1-x^2)^{m/2}\dfrac{d^m}{dx^m}P_l$ 取，不带上标负号；对照其他书时先确认这一条。}

\clearpage

\subhdr{目录（页码自动解析）}
\tabtop
\begin{tabular}{@{}p{0.10\linewidth}p{0.74\linewidth}r@{}}
\thead{节} & \thead{内容} & \thead{页} \\
\tabmidrule
\S0 & 本节要解决什么：主线与记号 & \pageref{sec:s0} \\
\S1 & 球坐标与它的 Laplace 算子 & \pageref{sec:s1} \\
\S2 & 第一次分离：径向与角向 & \pageref{sec:s2} \\
\S3 & 第二次分离：$\varphi$ 与 $\theta$ & \pageref{sec:s3} \\
\S4 & Legendre 方程、自然边界与球谐函数 & \pageref{sec:s4} \\
\S5 & 径向方程：Euler 方程与两组解 & \pageref{sec:s5} \\
\S6 & 球坐标 Laplace 的通解与定系数 & \pageref{sec:s6} \\
\S7 & 波动与输运：球贝塞尔与三重级数 & \pageref{sec:s7} \\
\S8 & 本征值问题总表（课件 p10 作业） & \pageref{sec:s8} \\
\S9 & 速查与易错 & \pageref{sec:s9} \\
\S10 & 三句话总结 & \pageref{sec:s10} \\
\S11 & 独立练习（五题） & \pageref{sec:practice} \\
附录 A & 五道练习的详解 & \pageref{sec:app} \\
\end{tabular}
\tabbot

\clearpage

% ===================== §0–§3 =====================
\hdr{\S0}{本节要解决什么：与前三节的关系}
\phantomsection\label{sec:s0}

§8.1 在矩形域上分离变量，§8.2 换到极坐标（圆、环），§8.3 加了第三个变量 $z$（柱）。这一节再换一次坐标：
球坐标 $(r,\theta,\varphi)$，区域是球、球外或球壳。

换坐标的理由和 §8.2 一样：\textbf{让边界落成坐标面}。球面 $r=r_0$ 在球坐标里就是「$r=$ 常数」，
和圆域里的 $r=\rho_0$ 是一回事。真正的麻烦不在 $r$，而在角向：球坐标的角向有\textbf{两个}变量
$\theta$ 与 $\varphi$，而 $\nabla^2$ 里的 $\sin\theta$、$\sin^2\theta$ 让 $\theta$ 的方程变成变系数方程。
这一节的全部新东西，都来自「怎么处理这个变系数的 $\theta$ 方程」。

方法骨架与 §8.1--§8.3 完全一样，只是多分离一次：

\begin{enumerate}[label=\textbf{(\arabic*)}]
  \item 第一次分离：$u=R(r)Y(\theta,\varphi)$，把 $r$ 与角向分开，得到一个常数 $\lambda$（\S2）；
  \item 第二次分离：$Y=\Theta(\theta)\Phi(\varphi)$，把 $\theta$ 与 $\varphi$ 分开，得到第二个常数 $\mu=m^2$（\S3）；
  \item 换元 $x=\cos\theta$，把 $\theta$ 方程变成\textbf{关联 Legendre 方程}；自然边界条件定出 $\lambda=l(l+1)$（\S4）；
  \item 径向方程在 $\lambda=l(l+1)$ 时是 Euler 方程，解为 $r^{l}$ 与 $r^{-(l+1)}$（\S5）；
  \item 球面条件定系数：用 $P_l^m$ 的正交性逐阶读出（\S6）；
  \item 若方程含时间，时间分离会多出一个 $k$，径向方程变成球贝塞尔方程（\S7）。
\end{enumerate}

要提醒一句：第 (3) 步的「自然边界条件」是本节的\textbf{主定理}，因为它同时干了两件事——
把 $\lambda$ 从连续谱钉成离散的 $l(l+1)$，并且决定了 $m$ 只能取到 $l$。不知道这一步的论证，
后面所有 $P_l^m$ 的求和上限都只是背下来的。

%=====================================================================
\hdr{\S1}{球坐标与它的 Laplace 算子}
\phantomsection\label{sec:s1}
%=====================================================================

\subhdr{1.1\quad 坐标、区域与三个退化点（课件 p1）}

球坐标与直角坐标的关系是
\[
x=r\sin\theta\cos\varphi,\qquad
y=r\sin\theta\sin\varphi,\qquad
z=r\cos\theta,
\]
其中 $r\ge0$、$0\le\theta\le\pi$、$0\le\varphi<2\pi$。

\begin{itemize}
  \item $r$：到原点的距离；球面就是 $r=r_0$。
  \item $\theta$：与 $+z$ 轴的夹角，叫极角（polar angle）。\textbf{不是}与 $xy$ 平面的夹角——这是最容易混的一处。
  \item $\varphi$：在 $xy$ 平面内自 $+x$ 轴起算的方位角，与 §8.2、§8.3 的 $\varphi$ 是同一个。
\end{itemize}

三个退化点先记下来，后面每一个都会用到：

\begin{enumerate}[label=(\arabic*)]
  \item \textbf{$r=0$}：原点处 $\theta,\varphi$ 都没有定义；球内问题的解必须在 $r\to0$ 时有界，这就是「自然边界条件」的第一处来源。
  \item \textbf{$\theta=0$ 与 $\theta=\pi$}：即 $+z$ 轴与 $-z$ 轴。这两条线上 $\varphi$ 无定义，所以解在 $\theta\to0,\pi$ 时也必须有界；
        换成 $x=\cos\theta$ 后，就是 $x=\pm1$ 有界。这是 $\lambda=l(l+1)$ 的来源（\S4）。
  \item \textbf{$\varphi=0$ 与 $\varphi=2\pi$}：是同一张半平面，所以 $\varphi$ 方向是周期条件，$\Phi(\varphi+2\pi)=\Phi(\varphi)$，给出 $\mu=m^2$（\S3）。
\end{enumerate}

一句话：球坐标里有\textbf{一个周期条件}（$\varphi$）与\textbf{两处自然边界}（$r=0$ 或 $r\to\infty$，以及 $x=\pm1$）。
三类条件各管一个常数，这就是本节要分离三次的原因。

\subhdr{1.2\quad 算子的来源：为什么长这样（课件 p1）}

课件 p1 直接写出球坐标下的 Laplace 算子：
\begin{equation}
\nabla^2u=\frac{1}{r^2}\frac{\partial}{\partial r}\Big(r^2\frac{\partial u}{\partial r}\Big)
+\frac{1}{r^2\sin\theta}\frac{\partial}{\partial\theta}\Big(\sin\theta\frac{\partial u}{\partial\theta}\Big)
+\frac{1}{r^2\sin^2\theta}\frac{\partial^2u}{\partial\varphi^2}.
\label{eq:lap}
\end{equation}

课件只写「Laplace operator in spherical coordinates」，没有给推导。这里补三件事，够用它，不展开九项链式法则。

\textbf{第一，三项各管什么。}第一项只含 $r$，是\textbf{径向}部分；第二、三项只含角向，是\textbf{角向}部分。
角向部分前面乘了 $1/r^2$，所以后面分离变量时可以把 $r$ 整体提出来——这就是第一次分离能成功的原因。

\textbf{第二，为什么第一项写成 $\frac1{r^2}\partial_r(r^2\partial_r\cdot)$ 而不是直接 $\partial_r^2$。}
在直角坐标里 $\partial_x^2+\partial_y^2+\partial_z^2$ 是「三个方向各求两次导」；换到球坐标后，
半径方向的「面积」随 $r^2$ 增长，散度的形式变了。这个写法等价于
\[
\frac{1}{r^2}\frac{d}{dr}\Big(r^2\frac{dR}{dr}\Big)=R''+\frac2rR',
\]
后面 \S5 解径向方程时会反复用它。两种写法都要认得。

\textbf{第三，用特例核对。}取球对称解 $u=u(r)$（与 $\theta,\varphi$ 无关），算子只剩第一项：
\[
\nabla^2u=\frac{1}{r^2}\frac{d}{dr}\Big(r^2\frac{du}{dr}\Big).
\]
取 $u=1/r$ 算一下：$r^2u'=-1$，再对 $r$ 求导得 $0$，故\allowbreak $\nabla^2(1/r)=0$（$r\ne0$）。

这正是静电学的点电荷电位，说明第一项没写错。再取 $u=r^2$：$r^2\cdot2r=2r^3$，求导得 $6r^2$，
除以 $r^2$ 得 $6$；而直角坐标下 $\nabla^2(x^2+y^2+z^2)=6$，两边一致。

\aside{如果要自己推一遍：从 $x,y,z$ 到 $r,\theta,\varphi$ 的链式法则共有九项，整理时需要用到
$\partial r/\partial x=x/r=\sin\theta\cos\varphi$ 这类关系。推导的细节不属于本讲范围；
只要记住三项的\textbf{结构}（径向一项、角向两项、公共因子 $1/r^2$），分离变量的每一步都不会错。}

%=====================================================================
\hdr{\S2}{第一次分离：径向与角向}
\phantomsection\label{sec:s2}
%=====================================================================

\subhdr{2.1\quad 设乘积解 $u=R(r)Y(\theta,\varphi)$（课件 p1）}

把 $u=R(r)Y(\theta,\varphi)$ 代入 \eqref{eq:lap}，两边除以 $RY$（先假定 $R,Y$ 不恒为零）：

\[
\frac{1}{R}\frac{d}{dr}\Big(r^2\frac{dR}{dr}\Big)
+\frac{1}{Y\sin\theta}\frac{\partial}{\partial\theta}\Big(\sin\theta\frac{\partial Y}{\partial\theta}\Big)
+\frac{1}{Y\sin^2\theta}\frac{\partial^2Y}{\partial\varphi^2}=0.
\]

第一项只含 $r$，后两项只含 $\theta,\varphi$。把后两项移到右边：
\begin{equation}
\underbrace{\frac{1}{R}\frac{d}{dr}\Big(r^2\frac{dR}{dr}\Big)}_{\text{只含 }r}
=\underbrace{-\frac{1}{Y\sin\theta}\frac{\partial}{\partial\theta}\Big(\sin\theta\frac{\partial Y}{\partial\theta}\Big)
-\frac{1}{Y\sin^2\theta}\frac{\partial^2Y}{\partial\varphi^2}}_{\text{只含 }\theta,\varphi}
\equiv\lambda .
\label{eq:sep1}
\end{equation}

\subhdr{2.2\quad 为什么两边必须同等于常数}

左边的式子里没有 $\theta,\varphi$，右边的式子里没有 $r$。可是两边恒等——对\textbf{任意}的 $r,\theta,\varphi$ 都成立。
固定一组 $(\theta,\varphi)$，让 $r$ 变：左边一般会变，右边却不变；既然两边永远相等，左边也只能不变。
同理右边也与 $\theta,\varphi$ 无关。于是两边都等于同一个常数，记作 $\lambda$。

这一步与 §8.2.2、§8.3.4 的论证是同一句话，只是变量换了。它的作用是把一个偏微分方程
\textbf{拆成两个}：一个只含 $r$，一个只含角向。

\subhdr{2.3\quad 拆出来的两个方程（课件 p2）}

由 \eqref{eq:sep1} 得

\textbf{径向方程}（Euler 型，课件 p2 左上）：
\begin{equation}
\frac{d}{dr}\Big(r^2\frac{dR}{dr}\Big)-\lambda R=0
\quad\Longleftrightarrow\quad
r^2R''+2rR'-\lambda R=0 .
\label{eq:radial}
\end{equation}

\textbf{角向方程}（球谐方程，课件 p2 右上）：
\begin{equation}
\frac{1}{\sin\theta}\frac{\partial}{\partial\theta}\Big(\sin\theta\frac{\partial Y}{\partial\theta}\Big)
+\frac{1}{\sin^2\theta}\frac{\partial^2Y}{\partial\varphi^2}+\lambda Y=0 .
\label{eq:angular}
\end{equation}

课件 p2 在角向方程旁边写了「$Y(\theta,\varphi)$: spherical harmonics」——球谐函数。这个名字先记住，
它的具体形状要等 \S4 把 $\lambda$ 定成 $l(l+1)$ 之后才写得出来。

\warnbox{\textbf{$\lambda$ 的符号还没定。}现在只知道它是一个常数，不知道正负。三种可能（$\lambda>0$、$=0$、$<0$）
都要在 \S4--\S5 里过一遍：$\lambda=l(l+1)\ge0$ 最后由自然边界条件逼出来，
而 $\lambda=0$ 是其中的 $l=0$ 特例（不是「平凡解」）。这与 §8.1 里「$\lambda=0$ 是否平凡要看边界」是同一件事。}

%=====================================================================
\hdr{\S3}{第二次分离：$\varphi$ 与 $\theta$}
\phantomsection\label{sec:s3}
%=====================================================================

\subhdr{3.1\quad 设 $Y=\Theta(\theta)\Phi(\varphi)$（课件 p2）}

角向方程 \eqref{eq:angular} 里还有两个变量。再设 $Y=\Theta(\theta)\Phi(\varphi)$，代入并乘以 $\sin^2\theta/(\Theta\Phi)$：

\[
\frac{\sin\theta}{\Theta}\frac{d}{d\theta}\Big(\sin\theta\frac{d\Theta}{d\theta}\Big)+\lambda\sin^2\theta
=-\frac{1}{\Phi}\frac{d^2\Phi}{d\varphi^2}.
\]

左边只含 $\theta$，右边只含 $\varphi$，故两边同等于常数，记作 $\mu$（课件 p2 中间）：
\begin{equation}
\frac{\sin\theta}{\Theta}\frac{d}{d\theta}\Big(\sin\theta\frac{d\Theta}{d\theta}\Big)+\lambda\sin^2\theta
=-\frac{1}{\Phi}\frac{d^2\Phi}{d\varphi^2}\equiv\mu .
\label{eq:sep2}
\end{equation}

拆成两个方程：

\textbf{$\varphi$ 的方程}：
\begin{equation}
\Phi''+\mu\Phi=0 .
\label{eq:phi}
\end{equation}

\textbf{$\theta$ 的方程}：
\begin{equation}
\sin\theta\frac{d}{d\theta}\Big(\sin\theta\frac{d\Theta}{d\theta}\Big)+(\lambda\sin^2\theta-\mu)\Theta=0 .
\label{eq:theta}
\end{equation}

\subhdr{3.2\quad 周期条件把 $\mu$ 钉成 $m^2$（课件 p2）}

$\varphi$ 方向的边界条件是周期条件：$\varphi$ 与 $\varphi+2\pi$ 是空间里同一条射线，所以
\[
\Phi(\varphi+2\pi)=\Phi(\varphi),\qquad \Phi'(\varphi+2\pi)=\Phi'(\varphi).
\]

方程 \eqref{eq:phi} 是常系数二阶方程，$\mu$ 的符号决定解的形状：

\begin{center}
\tabtop
\begin{tabular}{@{}L{0.10\linewidth}L{0.30\linewidth}L{0.48\linewidth}@{}}
\thead{$\mu$} & \thead{通解} & \thead{能否满足周期条件} \\
\tabmidrule
$\mu<0$ & $Ae^{\sqrt{-\mu}\,\varphi}+Be^{-\sqrt{-\mu}\,\varphi}$ & 不能：指数函数以 $2\pi$ 为周期只有恒为零 \\
$\mu=0$ & $A+B\varphi$ & 不能（除非 $B=0$，此时是常数，对应 $m=0$） \\
$\mu>0$ & $A\cos\sqrt{\mu}\,\varphi+B\sin\sqrt{\mu}\,\varphi$ & 能，但要求 $\sqrt{\mu}$ 是整数 \\
\end{tabular}
\tabbot
\end{center}

于是 $\mu=m^2$，$m=0,1,2,\dots$，特征函数为
\begin{equation}
\Phi_m(\varphi)=A_m\cos m\varphi+B_m\sin m\varphi .
\label{eq:phim}
\end{equation}

$m=0$ 时 $\Phi_0=A_0$ 是常数，对应轴对称（解与 $\varphi$ 无关）。这就是课件 p2 右侧那一支：
$\mu=m^2$ 与 $\Phi_m=A_m\cos m\varphi+B_m\sin m\varphi$。

\subhdr{3.3\quad $\theta$ 方程：本节真正的新东西（课件 p3 上半）}

把 $\mu=m^2$ 代回 \eqref{eq:theta}：
\begin{equation}
\sin\theta\frac{d}{d\theta}\Big(\sin\theta\frac{d\Theta}{d\theta}\Big)+(\lambda\sin^2\theta-m^2)\Theta=0 .
\label{eq:theta2}
\end{equation}

这个方程\textbf{没有初等解}——$\sin\theta$ 在分母、系数随 $\theta$ 变，直接猜解猜不出来。
课件 p3 的处理是换元 $x=\cos\theta$，把它变成 Legendre 方程。这一步是本节的分水岭，
\S4 专门做。

\aside{先把「为什么值得换元」说清：\eqref{eq:theta2} 在 $\theta=0,\pi$（即 $x=\pm1$）附近是奇异的，
而物理上解在这两点必须有界。$x=\cos\theta$ 把 $0\le\theta\le\pi$ 拉直成 $-1\le x\le1$，
奇异点就变成区间的两个端点；「解在端点有界」这件事，在 $x$ 的方程里变成一句标准的\textbf{自然边界条件}，
和 §8.2 圆内问题里「$r\to0$ 有界」是同一类条件。}

\subhdr{3.4\quad 与柱坐标的对照：为什么球坐标多分离了一次}

柱坐标 $(r,\varphi,z)$ 的 Laplace 只分离两次：$u=R(r)\Phi(\varphi)Z(z)$，因为 $z$ 方向是直的，
$Z''-\mu Z=0$ 直接解。球坐标的「第三个方向」$\theta$ 不是直线，而是一段角度，所以：

\begin{center}
\tabtop
\begin{tabular}{@{}L{0.14\linewidth}L{0.26\linewidth}L{0.24\linewidth}L{0.26\linewidth}@{}}
\thead{坐标} & \thead{第一次分离} & \thead{角向} & \thead{第三个方向} \\
\tabmidrule
直角 $(x,y,z)$ & $X(x)Y(y)Z(z)$ & — & $Z''-\lambda Z=0$，三角 \\
极坐标 $(r,\varphi)$ & $R(r)\Phi(\varphi)$ & $\Phi''+m^2\Phi=0$ & — \\
柱坐标 $(r,\varphi,z)$ & $R(r)\Phi(\varphi)Z(z)$ & $\Phi''+m^2\Phi=0$ & $Z''-\mu Z=0$，三角或双曲 \\
球坐标 $(r,\theta,\varphi)$ & $R(r)Y(\theta,\varphi)$ & $\Phi''+m^2\Phi=0$ & $\theta$：关联 Legendre 方程 \\
\end{tabular}
\tabbot
\end{center}

表中最后一列是本节唯一的「新方程」。其余部分——周期条件、Euler 方程、级数叠加、正交定系数——
与 §8.1--§8.3 逐条对应，只是换了特征函数。

\clearpage

% ===================== §4 Legendre =====================
\hdr{\S4}{Legendre 方程、自然边界与球谐函数}
\phantomsection\label{sec:s4}

这一节做三件事：把 $\theta$ 方程换成 $x=\cos\theta$ 的方程（\S4.1）；用 $x=\pm1$ 的有界性把 $\lambda$ 钉成 $l(l+1)$（\S4.2）；
写出解 $P_l$、$P_l^m$ 与它们的正交性（\S4.3--\S4.6）。课件 p3--p4 把这三件事压成两页，
其中「有界性为什么只留整数 $l$」与「$m$ 为什么只能取到 $l$」两处只给了结论，本节补上论证。

\subhdr{4.1\quad 换元 $x=\cos\theta$：从 $\theta$ 方程到 Legendre 方程（课件 p3）}

令
\[
x=\cos\theta,\qquad y(x)=\Theta(\theta).
\]
因为 $\theta\in[0,\pi]$，所以 $x\in[-1,1]$，且
\[
\frac{d}{d\theta}=\frac{dx}{d\theta}\frac{d}{dx}=-\sin\theta\frac{d}{dx},
\qquad \sin^2\theta=1-x^2 .
\]

把 $\dfrac{d\Theta}{d\theta}=-\sin\theta\,y'$ 代入 $\theta$ 方程 \eqref{eq:theta2}。先算第一项：
\[
\sin\theta\frac{d}{d\theta}\Big(\sin\theta\frac{d\Theta}{d\theta}\Big)
=\sin\theta\frac{d}{d\theta}\Big(-\sin^2\theta\,y'\Big)
=\sin\theta\cdot(-\sin\theta)\frac{d}{dx}\Big(-(1-x^2)y'\Big)
=\sin^2\theta\frac{d}{dx}\Big[(1-x^2)y'\Big].
\]
（第二步用了 $\frac{d}{d\theta}=-\sin\theta\frac{d}{dx}$；第三步把两个负号抵消。）

代回 \eqref{eq:theta2}，两边除以 $\sin^2\theta=1-x^2$：
\begin{equation}
\frac{d}{dx}\Big[(1-x^2)\frac{dy}{dx}\Big]+\Big(\lambda-\frac{m^2}{1-x^2}\Big)y=0 .
\label{eq:assoc}
\end{equation}

这就是课件 p3 的\textbf{关联 Legendre 方程}（Associated Legendre's equation）。$m=0$ 时它退化为
\begin{equation}
\frac{d}{dx}\Big[(1-x^2)\frac{dy}{dx}\Big]+\lambda y=0 ,
\label{eq:legendre}
\end{equation}
叫 \textbf{Legendre 方程}（课件 p4 下半）。

\warnbox{\textbf{换元时最容易丢的两处。}①$d/d\theta=-\sin\theta\,d/dx$，负号不能丢，但它会在下一步与另一个负号抵消，
所以最终方程里看不到；②$\sin^2\theta=1-x^2$，不是 $x^2-1$。这两处写错，得到的方程会差一个整体符号。
自检办法：$x=0$（$\theta=\pi/2$）附近，$(1-x^2)\approx1$，方程应近似 $y''+\lambda y\approx0$——这与 $\theta$ 方程的近似一致。}

\subhdr{4.2\quad 自然边界条件：为什么 $\lambda=l(l+1)$（课件 p3 下半）}

课件 p3 写的是：

\begin{quote}
\textbf{Eigenvalue $\lambda$?} $\to$ \textbf{Natural boundary condition}\\
The solution is finite at $x=1$ and $x=-1$ (i.e., $\theta=0$ and $\theta=\pi$), then we get
$\lambda=l(l+1)$, $l=0,1,2,3,\dots$ \quad (Chapter 12)
\end{quote}

「解在 $x=\pm1$ 有界」为什么就能把 $\lambda$ 钉成 $l(l+1)$？下面给完整的思路。结论先写出来：

\thm{定理 1（Legendre 方程的有界谱）}{方程 \eqref{eq:legendre} 在 $[-1,1]$ 上有非零有界解的充要条件是
\[
\lambda=l(l+1),\qquad l=0,1,2,\dots
\]
此时有界解就是 $l$ 次 Legendre 多项式 $P_l(x)$（差一个常数倍）。}

\pfline

\textbf{第一步：在 $x=0$ 附近求幂级数解。}设 $y=\sum_{k\ge0}a_kx^k$（$x=0$ 是方程的常点，幂级数解一定存在，收敛半径至少到最近的奇点 $x=\pm1$）。代入 \eqref{eq:legendre}：

\[
(1-x^2)\sum_k k(k-1)a_kx^{k-2}-2x\sum_k ka_kx^{k-1}+\lambda\sum_ka_kx^k=0 .
\]

把三项都整理成 $x^k$ 的系数，得到递推关系
\begin{equation}
a_{k+2}=\frac{k(k+1)-\lambda}{(k+2)(k+1)}\,a_k,\qquad k=0,1,2,\dots
\label{eq:recur}
\end{equation}
（偶数项与奇数项各自成链，$a_0,a_1$ 自由，所以有两个线性无关解。）

\textbf{第二步：截断的条件。}若 $\lambda=l(l+1)$（$l$ 为非负整数），则 $k=l$ 时分子
$l(l+1)-\lambda=0$，于是 $a_{l+2}=0$，此后所有同奇偶的系数都为零——级数\textbf{截断}成 $l$ 次多项式。
这就是 $P_l(x)$。反过来，若 $\lambda$ 不是这种形式，递推永远不会截断。

\textbf{第三步：不截断时在端点发散。}设 $\lambda$ 不是 $l(l+1)$。由 \eqref{eq:recur}，
当 $k$ 很大时 $a_{k+2}/a_k\to1$，级数的通项不趋于零（更精确地说，$a_k$ 只按 $1/k$ 的幂衰减，不够快），
在 $x=\pm1$ 处级数发散。具体地：把级数在 $x=1$ 附近重排，可以证明两个线性无关解中至少有一个在 $x=1$ 或 $x=-1$ 呈
\[
y\sim \ln(1-x)\quad\text{或}\quad y\sim\frac{1}{1-x}
\]
型的发散（哪一端发散取决于 $\lambda$ 落在哪两个整数之间）。总之：\textbf{只要不截断，就至少有一端无界}。

\textbf{第四步：结论。}物理上 $\theta=0,\pi$ 是两个极点，解在那里必须有限（否则场强发散、温度无穷）。
排除掉发散的解之后，只剩下截断的那一支，于是 $\lambda=l(l+1)$。\qedsym

\aside{证明里只用到一个事实：\textbf{两个线性无关解，一个在端点有界、另一个一般无界}。这与 §8.2 圆内问题
「$r=0$ 有界丢掉 $\ln r$ 与 $r^{-m}$」是同一件事——都是「自然边界条件筛掉奇异解」。
区别只在：圆内是一个端点（$r=0$），球面是两个端点（$x=\pm1$），所以球坐标的谱比极坐标多一个整数 $l$。}

\subhdr{4.3\quad Legendre 多项式：定义、前几个与 Rodrigues 公式（课件 p4）}

$\lambda=l(l+1)$ 时的有界解记作 $P_l(x)$，叫 $l$ 次 Legendre 多项式。它的标准定义（Rodrigues 公式）是
\begin{equation}
P_l(x)=\frac{1}{2^l\,l!}\frac{d^l}{dx^l}(x^2-1)^l,\qquad l=0,1,2,\dots
\label{eq:rodrigues}
\end{equation}

前几个（必须记住，后面定系数要用）：
\begin{align}
P_0(x)&=1, \label{eq:P0}\\
P_1(x)&=x, \label{eq:P1}\\
P_2(x)&=\frac{3x^2-1}{2}, \label{eq:P2}\\
P_3(x)&=\frac{5x^3-3x}{2}, \label{eq:P3}\\
P_4(x)&=\frac{35x^4-30x^2+3}{8}. \label{eq:P4}
\end{align}

三条常用性质：

\begin{enumerate}[label=\textbf{(\arabic*)}]
  \item \textbf{端点值}：$P_l(1)=1$，$P_l(-1)=(-1)^l$。代入 \eqref{eq:P0}--\eqref{eq:P4} 可验。
  \item \textbf{奇偶性}：$P_l(-x)=(-1)^lP_l(x)$。$l$ 偶是偶函数，$l$ 奇是奇函数——这直接决定球面对称的初值只激发哪些 $l$。
  \item \textbf{递推}：$(l+1)P_{l+1}=(2l+1)xP_l-lP_{l-1}$，用来从低阶算高阶，比每次求导快。
\end{enumerate}

\aside{\textbf{用 $P_l$ 表示「球面上的 $z$ 的多项式」。}因为 $x=\cos\theta=z/r$，所以
\[
P_0=1,\qquad P_1=\frac{z}{r},\qquad P_2=\frac{3z^2-r^2}{2r^2},\qquad P_3=\frac{5z^3-3zr^2}{2r^3}.
\]
在球面 $r=r_0$ 上，它们就是 $1$、$z/r_0$、$(3z^2-r_0^2)/(2r_0^2)$……——球面上的「多项式基」。
后面 \S6 的定系数，本质上就是把边界值按这组基展开。}

\subhdr{4.4\quad 正交性与模：$\int_{-1}^{1}P_lP_{l'}\,dx=\frac{2}{2l+1}\delta_{ll'}$}

Legendre 方程可以写成 Sturm--Liouville 形式
\[
\frac{d}{dx}\Big[(1-x^2)y'\Big]+\lambda y=0 ,
\]
其中 $p(x)=1-x^2$ 在端点为零、$q=0$、权函数 $w=1$。于是不同本征值的特征函数正交：
\begin{equation}
\int_{-1}^{1}P_l(x)P_{l'}(x)\,dx=\frac{2}{2l+1}\,\delta_{ll'} .
\label{eq:orthP}
\end{equation}

\textbf{推导（两行）。}把 $P_l$ 与 $P_{l'}$ 的方程分别乘以对方再相减：
\[
\frac{d}{dx}\Big[(1-x^2)(P_l'P_{l'}-P_lP_{l'}')\Big]+\big[l(l+1)-l'(l'+1)\big]P_lP_{l'}=0 .
\]
在 $[-1,1]$ 上积分，第一项是全导数，端点处 $1-x^2=0$，所以它为零；剩下的就是
$\big[l(l+1)-l'(l+1)\big]\int P_lP_{l'}=0$。$l\ne l'$ 时前因子不为零，故积分为零。模方 $2/(2l+1)$
可由 $P_l(1)=1$ 与递推关系算出，这里直接引用。

\textbf{三个模方值}（常用）：
\[
\int_{-1}^1P_0^2\,dx=2,\qquad
\int_{-1}^1P_1^2\,dx=\frac23,\qquad
\int_{-1}^1P_2^2\,dx=\frac25 .
\]

\warnbox{\textbf{权函数是 $1$，区间是 $[-1,1]$。}换成 $\theta$ 的积分时要乘 $\sin\theta\,d\theta$：
\[
\int_0^\pi P_l(\cos\theta)P_{l'}(\cos\theta)\sin\theta\,d\theta=\frac{2}{2l+1}\delta_{ll'} .
\]
这个式子后面定系数时直接用；漏掉 $\sin\theta$ 是最常见的错。}

\subhdr{4.5\quad 关联 Legendre 函数：$m\ne0$ 的那一支（课件 p4）}

回到 $m\ne0$ 的关联 Legendre 方程 \eqref{eq:assoc}。它的有界解是
\begin{equation}
P_l^m(x)=(1-x^2)^{m/2}\frac{d^m}{dx^m}P_l(x),\qquad 0\le m\le l .
\label{eq:assocP}
\end{equation}

三点说明：

\begin{enumerate}[label=\textbf{(\arabic*)}]
  \item \textbf{$m$ 只能取到 $l$。}$P_l$ 是 $l$ 次多项式，求 $l+1$ 次导数为零；所以 $m>l$ 时 $P_l^m\equiv0$，
        不是「另一个解」，而是没有解。这就是课件 p4 那行「$m=0,1,2,\dots,l-1,l$，Be careful about $m$!!!」的意思。
  \item \textbf{$m=0$ 退化}：$P_l^0=P_l$，即普通的 Legendre 多项式（课件 p4 下半：$m=0$ 时是轴对称问题，方程退化为 Legendre 方程）。
  \item \textbf{因子 $(1-x^2)^{m/2}=\sin^m\theta$}：它在 $\theta=0,\pi$ 处为零，保证 $P_l^m$ 在极点仍然有界
        （虽然导数可能发散，但函数值有限）。这个因子正是「$m$ 越大、角向越集中在赤道附近」的来源。
\end{enumerate}

前几个关联函数（$m=1,2$）：
\begin{align}
P_1^1(x)&=(1-x^2)^{1/2}=\sin\theta, \label{eq:P11}\\
P_2^1(x)&=3x(1-x^2)^{1/2}=3\cos\theta\sin\theta, \label{eq:P21}\\
P_2^2(x)&=3(1-x^2)=3\sin^2\theta . \label{eq:P22}
\end{align}

正交性（固定 $m$，对 $l$ 正交）：
\begin{equation}
\int_{-1}^{1}P_l^m(x)P_{l'}^m(x)\,dx=\frac{2}{2l+1}\cdot\frac{(l+m)!}{(l-m)!}\,\delta_{ll'} .
\label{eq:orthPm}
\end{equation}

注意模方里多出的因子 $(l+m)!/(l-m)!$：$m$ 越大，模方越大，所以同样大小的边界值在 $m$ 大的模态上
系数要更小。这一条在 \S6 定系数时会直接用。

\subhdr{4.6\quad 球谐函数：把 $\theta$ 与 $\varphi$ 合起来}

把 \eqref{eq:phim} 与 \eqref{eq:assocP} 乘起来，得到角向方程 \eqref{eq:angular} 的解。实形式的球谐函数是
\begin{equation}
Y_{lm}(\theta,\varphi)=P_l^m(\cos\theta)\cos m\varphi
\quad\text{与}\quad
\tilde Y_{lm}(\theta,\varphi)=P_l^m(\cos\theta)\sin m\varphi ,
\qquad 0\le m\le l .
\label{eq:spherical}
\end{equation}
（$m=0$ 时 $\sin0\varphi\equiv0$，只剩 $P_l(\cos\theta)$ 一支。）

它们满足角向方程
\[
\frac{1}{\sin\theta}\partial_\theta(\sin\theta\,\partial_\theta Y)+\frac{1}{\sin^2\theta}\partial_\varphi^2Y+l(l+1)Y=0 ,
\]
并且在球面上正交：
\begin{equation}
\int_0^{2\pi}\!\!\int_0^{\pi}Y_{lm}Y_{l'm'}\,\sin\theta\,d\theta\,d\varphi
=\frac{2\pi}{2l+1}\frac{(l+m)!}{(l-m)!}\,\delta_{ll'}\delta_{mm'}\cdot
\begin{cases}
2,& m=0,\\
1,& m\ge1,
\end{cases}
\label{eq:orthY}
\end{equation}
其中 $\sin m\varphi$ 与 $\cos m\varphi$ 之间也正交（对 $\varphi$ 积分为零）。

\aside{\textbf{为什么只取 $m\ge0$？}负的 $m$ 并不给出新函数：$\cos(-m\varphi)=\cos m\varphi$、$\sin(-m\varphi)=-\sin m\varphi$，
只是同一个函数差一个符号。所以在实形式里把 $m$ 限制在 $0\le m\le l$，一共 $2l+1$ 个独立函数（$m=0$ 一个，
$m=1,\dots,l$ 各两个）。这就是「$l$ 层有 $2l+1$ 个球谐」的来源。}

\subhdr{4.7\quad 教学小例：最低三阶球谐在球面上的样子}

取球面 $r=r_0$，把 $P_l$ 写成直角坐标（用 $x=\cos\theta=z/r_0$）：

\begin{center}
\tabtop
\begin{tabular}{@{}L{0.10\linewidth}L{0.30\linewidth}L{0.50\linewidth}@{}}
\thead{$l$} & \thead{$P_l(\cos\theta)$} & \thead{球面上的形状} \\
\tabmidrule
$0$ & $1$ & 常数：整个球面同值 \\
$1$ & $\cos\theta=z/r_0$ & 北半球正、南半球负，赤道为零：一个偶极分布 \\
$2$ & $\dfrac{3\cos^2\theta-1}{2}=\dfrac{3z^2-r_0^2}{2r_0^2}$ & 两极正、赤道负（或反过来）：一个四极分布 \\
\end{tabular}
\tabbot
\end{center}

\aside{这三行可以直接用来核对后面的解。例如球内 Laplace 问题，边界值 $u(r_0,\theta)=U_0\cos\theta$，
只含 $l=1$ 的 $P_1$，所以解只能是 $u=C r\cos\theta$；在 $r=r_0$ 上定出 $C r_0=U_0$，
于是 $u=U_0\,(r/r_0)\cos\theta=U_0 z/r_0$——一个沿 $z$ 方向的均匀场。这与 §8.2 里「边界值只含一个谐波，
解就只含那一项」是同一句话，只是基函数从 $\cos m\varphi$ 换成了 $P_l(\cos\theta)$。}

\clearpage

% ===================== §5–§7 =====================
\hdr{\S5}{径向方程：Euler 方程与两组解}
\phantomsection\label{sec:s5}

\subhdr{5.1\quad $\lambda=l(l+1)$ 时的径向方程（课件 p5）}

把 $\lambda=l(l+1)$ 代回径向方程 \eqref{eq:radial}：
\begin{equation}
\frac{d}{dr}\Big(r^2\frac{dR}{dr}\Big)-l(l+1)R=0
\quad\Longleftrightarrow\quad
r^2R''+2rR'-l(l+1)R=0 .
\label{eq:euler}
\end{equation}

这是\textbf{Euler 方程}（变系数，但每项的次数与求导阶数相同）。试解 $R=r^\alpha$：$R'=\alpha r^{\alpha-1}$，$R''=\alpha(\alpha-1)r^{\alpha-2}$，代入得
\[
\alpha(\alpha-1)+2\alpha-l(l+1)=\alpha^2+\alpha-l(l+1)=0 .
\]
即
\[
(\alpha-l)\big(\alpha+(l+1)\big)=0
\quad\Longrightarrow\quad
\alpha=l\ \ \text{或}\ \ \alpha=-(l+1).
\]

所以
\begin{equation}
R_l(r)=C_l\,r^{l}+D_l\,r^{-(l+1)} .
\label{eq:Rl}
\end{equation}

$l=0$ 时两项是 $1$ 与 $1/r$——这正是静电学的常数电位与点电荷电位（\S1.2 用 $u=1/r$ 核对过算子）。

\aside{\textbf{与 §8.3 柱坐标 Euler 方程的区别。}柱坐标的径向方程是
\[
r^2R''+rR'-m^2R=0\quad\Rightarrow\quad R=C r^m+D r^{-m},
\]
指数是 $\pm m$；球坐标的方程多了一个 $r$（$2rR'$ 而不是 $rR'$），指数变成 $l$ 与 $-(l+1)$，两者之和为 $-1$。
这个差别来自球坐标的「面积」按 $r^2$ 增长、柱坐标按 $r$ 增长。记的时候把两条放在一起看，不容易混。}

\subhdr{5.2\quad 三种区域，各丢哪一项（课件 p5--p6）}

$r^l$ 与 $r^{-(l+1)}$ 在 $r=0$ 与 $r\to\infty$ 的行为相反，所以「留哪一项」完全由区域决定：

\begin{center}
\tabtop
\begin{tabular}{@{}L{0.16\linewidth}L{0.25\linewidth}L{0.25\linewidth}L{0.24\linewidth}@{}}
\thead{区域} & \thead{$r^l$} & \thead{$r^{-(l+1)}$} & \thead{保留} \\
\tabmidrule
球内 $0\le r<r_0$ & $r\to0$ 时有界 & $r\to0$ 时发散 & 只留 $r^l$ \\
球外 $r>r_0$ & $r\to\infty$ 时发散 & 有界（趋于零） & 只留 $r^{-(l+1)}$ \\
球壳 $a<r<b$ & 两端都有界 & 两端都有界 & 两组都留 \\
\end{tabular}
\tabbot
\end{center}

\textbf{三个例外必须记住：}

\begin{enumerate}[label=\textbf{(\arabic*)}]
  \item \textbf{球外如果还有外场（如均匀电场 $E_0$），$r^l$ 不能丢。}因为远处不是零，而是给定的外场；
        外场本身在球坐标里就是 $r^lP_l$ 型的项（$E_0r\cos\theta=E_0rP_1$）。此时通解两组都留，
        用「$r\to\infty$ 时趋于给定外场」与「球面条件」各定一部分系数。§8.2 的导线例题就是这个结构。
  \item \textbf{球内如果有源（点电荷、点热源），$r^{-(l+1)}$ 不能丢。}源在原点时，解在 $r=0$ 本来就不有界；
        这时的「自然边界条件」不再是「有界」，而是「在源以外有界、在源处给定强度」。
  \item \textbf{$l=0$ 的 $r^{-(l+1)}=1/r$ 是球对称的点源解。}它在 $r=0$ 发散，但在球外问题里是合法的。
\end{enumerate}

\subhdr{5.3\quad 球内问题为什么只留 $r^l$：一个直观的核对}

取球内 $l=1$ 的两项：$r$ 与 $r^{-2}$。前者在原点为零，后者在原点无穷大。
若边界条件给的是球面 $u(r_0,\theta)=U_0\cos\theta$（有限），那么解在球内也必须有限（否则球心处物理量无意义），
所以只能留 $r$。\S6 会算出 $u=U_0(r/r_0)\cos\theta$，在 $r=0$ 处为零，符合。

%=====================================================================
\hdr{\S6}{球坐标 Laplace 的通解与定系数}
\phantomsection\label{sec:s6}
%=====================================================================

\subhdr{6.1\quad 通解（课件 p5）}

把径向 \eqref{eq:Rl} 与角向 \eqref{eq:spherical} 乘起来并叠加，得到球坐标 Laplace 方程的通解：
\begin{equation}
\begin{aligned}
u(r,\theta,\varphi)={}&\sum_{l=0}^{\infty}\sum_{m=0}^{l}
\big(A_l^m r^l+C_l^m r^{-(l+1)}\big)P_l^m(\cos\theta)\cos m\varphi\\
&+\sum_{l=0}^{\infty}\sum_{m=0}^{l}
\big(B_l^m r^l+D_l^m r^{-(l+1)}\big)P_l^m(\cos\theta)\sin m\varphi .
\end{aligned}
\label{eq:general}
\end{equation}

这就是课件 p5 那个双和：一组 $\cos m\varphi$、一组 $\sin m\varphi$，每组里 $r^l$ 与 $r^{-(l+1)}$ 各带两个系数。

\textbf{读法}：对每一个 $(l,m)$，径向有两个自由度（$r^l$、$r^{-(l+1)}$），角向也有两个（$\cos m\varphi$、$\sin m\varphi$），
所以一对 $(l,m)$ 最多四个待定系数。区域一确定（丢哪一项），自由度立减一半。

\subhdr{6.2\quad 球内：只留 $r^l$（课件 p5 下半）}

球内 $r<r_0$，丢掉 $r^{-(l+1)}$：
\begin{equation}
u(r,\theta,\varphi)=\sum_{l=0}^{\infty}\sum_{m=0}^{l}
r^l\big(a_l^m\cos m\varphi+b_l^m\sin m\varphi\big)P_l^m(\cos\theta) .
\label{eq:inside}
\end{equation}

轴对称（$m=0$）时只剩
\begin{equation}
u(r,\theta)=\sum_{l=0}^{\infty}C_l\,r^lP_l(\cos\theta) .
\label{eq:axisym}
\end{equation}

这就是课件 p6 上半的式子。

\subhdr{6.3\quad 球外：课件留的那个问题（课件 p5 下）}

课件 p5 在球内公式下面写了一个红框：\textbf{「The solution outside a sphere?」}——没有给答案。补在这里。

球外 $r>r_0$，若远处没有外场（$u\to0$ 当 $r\to\infty$），丢掉 $r^l$：
\begin{equation}
u(r,\theta,\varphi)=\sum_{l=0}^{\infty}\sum_{m=0}^{l}
r^{-(l+1)}\big(a_l^m\cos m\varphi+b_l^m\sin m\varphi\big)P_l^m(\cos\theta) .
\label{eq:outside}
\end{equation}

\textbf{若球外还有均匀外场}（例如导体球放在均匀电场 $E_0\hat z$ 里），则远处 $u\to-E_0z=-E_0r\cos\theta=-E_0rP_1(\cos\theta)$，
这是一个 $l=1$ 的 $r^l$ 项，不能丢。此时解写成
\[
u=-E_0r\cos\theta+\sum_{l,m}r^{-(l+1)}(\cdots)P_l^m(\cos\theta)(\cdots),
\]
球面 $u(r_0)=0$ 再把 $l=1$ 的 $r^{-2}$ 系数定出来，最终
\[
u=-E_0r\cos\theta+E_0\frac{r_0^3}{r^2}\cos\theta
=-E_0\Big(r-\frac{r_0^3}{r^2}\Big)\cos\theta .
\]
（这个结果与 §8.2 的导线例题同构：都是「均匀外场 + 一个球/柱的感应项」，只是柱坐标的感应项是 $r^{-1}$、球坐标是 $r^{-2}$。）

\subhdr{6.4\quad 球壳：两组都留}

球壳 $a<r<b$：$r^l$ 与 $r^{-(l+1)}$ 在 $[a,b]$ 上都有界，所以两组都留：
\[
u=\sum_{l,m}\big(A_l^mr^l+C_l^mr^{-(l+1)}\big)P_l^m(\cos\theta)\cos m\varphi+\cdots .
\]
内球面与外球面各给一个条件，正好定下每一阶的两个径向系数——与 §8.2 环域问题的结构完全一样。

\subhdr{6.5\quad 定系数：用 $P_l^m$ 的正交性逐阶读出（课件 p6）}

设球面条件为
\[
u(r_0,\theta,\varphi)=f(\theta,\varphi).
\]
把 \eqref{eq:inside} 在 $r=r_0$ 取值，得
\[
f(\theta,\varphi)=\sum_{l,m}r_0^l\big(a_l^m\cos m\varphi+b_l^m\sin m\varphi\big)P_l^m(\cos\theta).
\]
两边乘 $P_{l'}^{m'}(\cos\theta)\cos m'\varphi$，再对 $\theta\in[0,\pi]$、$\varphi\in[0,2\pi]$ 积分（权 $\sin\theta\,d\theta\,d\varphi$），
用正交性 \eqref{eq:orthY} 只剩 $(l,m)=(l',m')$ 那一项，于是
\begin{equation}
a_l^m=\frac{1}{r_0^l}\cdot\frac{(2l+1)}{2\pi}\cdot\frac{(l-m)!}{(l+m)!}\cdot\frac{1}{\epsilon_m}
\int_0^{2\pi}\!\!\int_0^{\pi}f(\theta,\varphi)P_l^m(\cos\theta)\cos m\varphi\,\sin\theta\,d\theta\,d\varphi ,
\label{eq:coef}
\end{equation}
其中 $\epsilon_0=2$、$\epsilon_m=1$（$m\ge1$）。$b_l^m$ 把 $\cos m\varphi$ 换成 $\sin m\varphi$ 即可。

这组公式与 \eqref{eq:orthPm} 的模方是同一件事，只是多除了 $r_0^l$。实际做题时，多数题目只含少数几阶，
不必套公式——直接比较系数更快。

\subhdr{6.6\quad 教学小例：三个能直接看出答案的球面条件}

\textbf{例 1（一阶）}：球内，$u(r_0,\theta)=U_0\cos\theta$。

边界只含 $P_1(\cos\theta)$，所以解只含 $l=1$ 的 $r$ 项：$u=C_1rP_1(\cos\theta)=C_1r\cos\theta$。
在 $r=r_0$ 上：$C_1r_0=U_0$，故
\[
u=U_0\frac{r}{r_0}\cos\theta=U_0\frac{z}{r_0}.
\]
这是一个沿 $z$ 方向的均匀场（在球内）。核对：$\nabla^2z=0$ $\checkmark$，球面上取值 $U_0z/r_0=U_0\cos\theta$ $\checkmark$。

\textbf{例 2（零阶）}：球内，$u(r_0,\theta)=U_0$（常数）。

边界只含 $P_0=1$，解只含 $l=0$ 的 $r^0$ 项：$u=C_0$，$C_0=U_0$。解在球内是常数。
核对：常数的 $\nabla^2$ 为零 $\checkmark$；这与「无内热源的稳态温度场在区域内部没有极值」一致——边界处处同温，内部只能同温。

\textbf{例 3（混合两阶）}：球内，$u(r_0,\theta)=U_0\cos^2\theta$。

先把边界值按 $P_l$ 展开。用 $x=\cos\theta$：
\[
x^2=\frac{2}{3}P_2(x)+\frac{1}{3}P_0(x)
\quad\Longrightarrow\quad
U_0\cos^2\theta=\frac{U_0}{3}P_0+\frac{2U_0}{3}P_2 .
\]
所以解为
\[
u=\frac{U_0}{3}+\frac{2U_0}{3}\Big(\frac{r}{r_0}\Big)^2P_2(\cos\theta)
=\frac{U_0}{3}+\frac{U_0}{3}\Big(\frac{r}{r_0}\Big)^2(3\cos^2\theta-1).
\]
核对：$r=r_0$ 时回到 $U_0\cos^2\theta$ $\checkmark$；$r=0$ 时为 $U_0/3$（有限） $\checkmark$。
（这一步用了 $x^2=\frac{2}{3}P_2+\frac13P_0$，可由 $P_0=1$、$P_2=(3x^2-1)/2$ 反解得到。）

\subhdr{6.7\quad 课件 p6 的两条 Remark}

课件 p6 底部写了两条 Remark，逐条解释：

\begin{enumerate}[label=\textbf{(\arabic*)}]
  \item \textbf{「本征值 $\mu$ 与 $\lambda$ 由周期边界条件与自然边界条件确定」}——
        $\mu=m^2$ 来自 $\varphi$ 的周期条件（\S3.2），$\lambda=l(l+1)$ 来自 $x=\pm1$ 的有界性（\S4.2）。
        这两条都\textbf{不}来自球面上的边界条件；球面条件只用来定系数。
  \item \textbf{「球面边界条件可以是非齐次的」}——这与 §8.1 矩形域要先做「齐次化」不同：球坐标里
        $r=r_0$ 上的非齐次数据直接作为定系数的入口，不需要先平移。原因与 §8.2 环域相同：
        非齐次数据落在 $r$ 方向，而 $r$ 方向本来就没有「谱」要定，所以不冲突。
\end{enumerate}

%=====================================================================
\hdr{\S7}{波动与输运：球贝塞尔与三重级数}
\phantomsection\label{sec:s7}
%=====================================================================

\subhdr{7.1\quad 时间分离：从波动方程到 Helmholtz（课件 p7）}

齐次波动方程 $u_{tt}-a^2\nabla^2u=0$。设 $u(\mathbf r,t)=T(t)V(\mathbf r)$，代入并除以 $a^2TV$：
\[
\frac{T''}{a^2T}=\frac{\nabla^2V}{V}\equiv-k^2 .
\]
（为什么取 $-k^2$：波动问题的解在时间上要振荡，$T''<0$，所以这个常数必须为负。若取 $+k^2$，时间解是指数，
与「振动」不符。这与 §8.3.1.1 的论证是同一句。）

于是得到
\begin{equation}
\nabla^2V+k^2V=0\quad\text{（Helmholtz 方程）},\qquad
T''+a^2k^2T=0 .
\label{eq:helm}
\end{equation}

\textbf{输运方程}走同一条路：$u_t-a^2\nabla^2u=0$，分离得
\[
\nabla^2V+k^2V=0,\qquad T'+a^2k^2T=0 .
\]
空间部分完全相同，只有时间方程差一阶——所以输运的时间因子是 $e^{-a^2k^2t}$，只会衰减（\S7.5）。

\subhdr{7.2\quad Helmholtz 在球坐标里的分离（课件 p7 下半）}

把 \eqref{eq:lap} 里的 $\nabla^2$ 换成 $\nabla^2+k^2$，再设 $V=R(r)\Theta(\theta)\Phi(\varphi)$。
与 Laplace 情形同样的两步分离给出：
\begin{equation}
\Phi''+m^2\Phi=0 ,
\label{eq:phi2}
\end{equation}
\begin{equation}
\sin\theta\frac{d}{d\theta}\Big(\sin\theta\frac{d\Theta}{d\theta}\Big)+\big[\lambda\sin^2\theta-m^2\big]\Theta=0 ,
\label{eq:theta3}
\end{equation}
\begin{equation}
\frac{d}{dr}\Big(r^2\frac{dR}{dr}\Big)+\big[r^2k^2-\lambda\big]R=0 .
\label{eq:radialk}
\end{equation}

注意前两个方程与 Laplace 情形\textbf{一字不差}（\eqref{eq:phi2} 同 \eqref{eq:phi}，\eqref{eq:theta3} 同 \eqref{eq:theta2}），
所以 $\mu=m^2$、$\lambda=l(l+1)$ 的结论原封不动。变化只在径向：多出的 $r^2k^2$ 让 Euler 方程变成球贝塞尔方程。

\subhdr{7.3\quad 球贝塞尔方程：两步换元（课件 p8）}

取 $\lambda=l(l+1)$，径向方程 \eqref{eq:radialk} 写成
\begin{equation}
R''+\frac{2}{r}R'+\Big(k^2-\frac{l(l+1)}{r^2}\Big)R=0 .
\label{eq:sphbessel}
\end{equation}
这就是\textbf{球贝塞尔方程}。

\textbf{第一步：令 $x=kr$。}则 $d/dr=k\,d/dx$，$d^2/dr^2=k^2d^2/dx^2$，方程变成
\begin{equation}
\frac{d^2R}{dx^2}+\frac{2}{x}\frac{dR}{dx}+\Big(1-\frac{l(l+1)}{x^2}\Big)R=0 .
\label{eq:sphbesselx}
\end{equation}

\textbf{第二步：令 $R=x^{-1/2}y(x)$。}这一步的目的是消去一阶导项。算一下：
\[
R'=-\tfrac12x^{-3/2}y+x^{-1/2}y',\qquad
R''=\tfrac34x^{-5/2}y-x^{-3/2}y'+x^{-1/2}y'' .
\]
代入 \eqref{eq:sphbesselx}，乘 $x^{1/2}$ 整理：
\begin{equation}
y''+\frac{1}{x}y'+\Big[1-\frac{(l+\frac12)^2}{x^2}\Big]y=0 .
\label{eq:halfbessel}
\end{equation}

这正是\textbf{半整数阶贝塞尔方程}（阶数 $\nu=l+\tfrac12$）。它的两个线性无关解是 $J_{l+1/2}(x)$ 与 $N_{l+1/2}(x)$，
所以
\begin{equation}
R_l(r)=C_l\,j_l(kr)+D_l\,n_l(kr),
\label{eq:Rbessel}
\end{equation}
其中
\begin{equation}
j_l(x)=\sqrt{\frac{\pi}{2x}}\,J_{l+1/2}(x),\qquad
n_l(x)=\sqrt{\frac{\pi}{2x}}\,N_{l+1/2}(x)
\label{eq:jlnl}
\end{equation}
叫第一类、第二类\textbf{球贝塞尔函数}（课件 p8 的 the first / second kind spherical Bessel functions）。

\aside{\textbf{为什么值得记 $j_0,n_0$。}低阶的球贝塞尔是初等函数：
\[
j_0(x)=\frac{\sin x}{x},\qquad n_0(x)=-\frac{\cos x}{x},
\]
\[
j_1(x)=\frac{\sin x}{x^2}-\frac{\cos x}{x},\qquad
n_1(x)=-\frac{\cos x}{x^2}-\frac{\sin x}{x}.
\]
球对称问题（$l=0$）只用 $j_0,n_0$，可以手算到底。例如球内 $u(r_0)=0$ 的球对称振动，
条件 $j_0(kr_0)=0$ 就是 $\sin(kr_0)=0$，得 $k=n\pi/r_0$——与一维弦的 $n\pi/\ell$ 同一个形式。}

\subhdr{7.4\quad $k$ 由球面齐次条件定（课件 p8 下半）}

课件 p8 写：\textbf{「Eigenvalue $k$ is determined by the homogeneous boundary condition of the surface of the sphere.」}

设球面齐次条件为 $u(r_0,\theta,\varphi)=0$（例如固定零温、理想导体球面）。球内解只含 $j_l$（$n_l$ 在 $r=0$ 发散，丢掉）：
\[
R_l(r)=C_l\,j_l(kr).
\]
球面条件 $j_l(kr_0)=0$ 要求 $kr_0$ 是 $j_l$ 的零点。记 $j_l$ 的第 $n$ 个正零点为 $x_{ln}$，则
\begin{equation}
k=k_{ln}=\frac{x_{ln}}{r_0},\qquad n=1,2,3,\dots
\label{eq:kln}
\end{equation}

几个常用的零点：$j_0$ 的零点是 $n\pi$（$x_{0n}=n\pi$）；$j_1$ 的零点是 $\tan x=x$ 的正根（$x_{11}\approx4.4934$，
$x_{12}\approx7.7253$）；$j_2$ 的零点约为 $5.7635,9.0950,\dots$。

\textbf{与柱坐标的对照}：柱坐标的 $k$ 由 $J_m(ka)=0$ 定（$k=x_n^{(m)}/a$），零点间隔不等；
球坐标的 $k$ 由 $j_l(kr_0)=0$ 定，$l=0$ 时零点等距（$n\pi$），$l\ge1$ 时间隔也不等。
共同点：\textbf{谱的离散性来自边界上的零点，而不是来自区间的「长度」}。

\subhdr{7.5\quad 三重级数：波动与输运的通解（课件 p9）}

\textbf{波动}：$T''+a^2k^2T=0$ 的解是 $\cos(kat)$ 与 $\sin(kat)$。把三部分乘起来叠加：
\begin{equation}
u(r,\theta,\varphi,t)=\sum_{n=1}^{\infty}\sum_{l=0}^{\infty}\sum_{m=0}^{l}
R_l(k_{ln}r)\,P_l^m(\cos\theta)\,\Phi_m(\varphi)\,
\big[C_{lnm}\cos(k_{ln}at)+D_{lnm}\sin(k_{ln}at)\big] .
\label{eq:wavesol}
\end{equation}

\textbf{输运}：$T'+a^2k^2T=0$ 的解是 $e^{-a^2k^2t}$：
\begin{equation}
u(r,\theta,\varphi,t)=\sum_{n=1}^{\infty}\sum_{l=0}^{\infty}\sum_{m=0}^{l}
R_l(k_{ln}r)\,P_l^m(\cos\theta)\,\Phi_m(\varphi)\,
E_{lnm}\,e^{-a^2k_{ln}^2t} .
\label{eq:transportsol}
\end{equation}

\textbf{三重和的指标来源}（三个指标各管一件事，缺一不可）：

\begin{center}
\tabtop
\begin{tabular}{@{}L{0.10\linewidth}L{0.22\linewidth}L{0.56\linewidth}@{}}
\thead{指标} & \thead{来自哪里} & \thead{由什么条件定} \\
\tabmidrule
$m$ & $\varphi$ 方向 & 周期条件 $\Phi(\varphi+2\pi)=\Phi(\varphi)$，$\mu=m^2$ \\
$l$ & $\theta$ 方向 & $x=\pm1$ 有界（自然边界），$\lambda=l(l+1)$ \\
$n$ & $r$ 方向 & 球面齐次条件 $j_l(kr_0)=0$，$k=k_{ln}$ \\
\end{tabular}
\tabbot
\end{center}

这与 §8.3 柱坐标的「三重和 $m,n,j$」是同一张表，只是第三个指标的来源从「$z$ 的两端条件」
换成了「球面的零点条件」。

\textbf{系数怎么定。}初始条件 $u(r,\theta,0)$ 与 $u_t(r,\theta,0)$ 分别在球内展开：
先用 $j_l(k_{ln}r)$ 在 $[0,r_0]$ 上（带权 $r^2$）的正交性定 $n$，再用 $P_l^m$ 与 $\cos m\varphi/\sin m\varphi$ 定 $l,m$。
三步正交性各管一个指标，顺序可以换，但不能漏。

%=====================================================================

% ===================== §8–§10 =====================
\hdr{\S8}{本征值问题总表（课件 p10 作业）}
\phantomsection\label{sec:s8}

课件 p10 的作业只有一句话：

\begin{quote}
\textbf{Assignments}\quad Summarize the eigenvalue problems of Laplace's equation, wave equation
and transport equation in polar, cylindrical and spherical coordinates.
\end{quote}

这句话要求把 §8.1--§8.4 的谱问题并成一张表。本节给三张表（Laplace 一张、波与输运各一张）与一段判据，
读者可以拿它当复习提纲。

\subhdr{8.1\quad Laplace 方程：三种坐标的谱}

\begin{center}
\tabtop
\begin{tabular}{@{}L{0.10\linewidth}L{0.17\linewidth}L{0.27\linewidth}L{0.26\linewidth}@{}}
\thead{坐标} & \thead{区域} & \thead{谱（由什么条件定）} & \thead{另一个方向} \\
\tabmidrule
极 & 圆、环 & $\lambda=m^2$（周期条件），$m=0,1,2,\dots$；特征函数 $1,\cos m\varphi,\sin m\varphi$ & $r^m$、$r^{-m}$（$m=0$：$1,\ln r$） \\
柱 & 圆柱（$z$ 有界） & 角向同上；$z$ 方向 $\mu$：哪组边界齐次定符号，$\mu>0$ 贝塞尔、$\mu<0$ 三角、$\mu=0$ 欧拉 & $J_m(kr)$ 或 $I_m$ 或 $r^{\pm m}$ \\
球 & 球、球壳 & $\lambda=l(l+1)$（$x=\pm1$ 有界），$m=0,\dots,l$；特征函数 $P_l^m(\cos\theta)\cos m\varphi$、$\sin m\varphi$ & $r^l$、$r^{-(l+1)}$ \\
\end{tabular}
\tabbot
\end{center}

\textbf{三行的共同点}：角向（$\varphi$）永远是周期条件给 $m^2$；差别在「第二个方向」——
极坐标的 $r$ 是 Euler、柱坐标的 $z$ 有三种情形、球坐标的 $\theta$ 是 Legendre。
三者的谱都是\textbf{离散}的，因为都有一组齐次边界/周期/自然条件。

\subhdr{8.2\quad 波动与输运：多出的时间因子}

\begin{center}
\tabtop
\begin{tabular}{@{}L{0.10\linewidth}L{0.24\linewidth}L{0.24\linewidth}L{0.26\linewidth}@{}}
\thead{坐标} & \thead{空间谱（同表 8-1）} & \thead{径向/第三方向（含 $k$）} & \thead{时间因子} \\
\tabmidrule
极 & $m^2$ & $J_m(kr)$，$k=x_n^{(m)}/a$ & 波：$\cos(kat),\sin(kat)$；输运：$e^{-a^2k^2t}$ \\
柱 & $m^2$ 与 $\mu$ 的三种符号 & $J_m$、$\cos/\sin$、$e^{\pm\nu z}$ 三者组合 & 同上 \\
球 & $l(l+1)$、$m^2$ & $j_l(kr)$，$k=x_{ln}/r_0$ & 同上 \\
\end{tabular}
\tabbot
\end{center}

\textbf{读法}：时间是最后分离的（或最先，取决于你从哪一步开始），它只多一个因子，不改变空间谱。
所以「波动与输运的本征值问题」=「Laplace 的空间谱」+「球面/边界上的 $k$ 条件」+「时间因子」。

\subhdr{8.3\quad 三个判据：本征值到底由谁定}

把 §8.1--§8.4 里所有「定本征值」的话压缩成三条：

\begin{enumerate}[label=\textbf{(\arabic*)}]
  \item \textbf{周期条件定角向}：$\Phi(\varphi+2\pi)=\Phi(\varphi)\Rightarrow\mu=m^2$。三种坐标共用。
  \item \textbf{自然边界条件定「退化点」}：$r=0$ 有界丢掉奇异解（极、柱、球的径向）；$x=\pm1$ 有界定 $\lambda=l(l+1)$（球的 $\theta$）。
  \item \textbf{齐次边界条件定「有限区间」}：矩形两端定 $n\pi/\ell$；圆/圆柱/球的表面定 $k$（$J_m$ 或 $j_l$ 的零点）。
\end{enumerate}

一句话：\textbf{谱来自「齐次」的那一处}。哪一组条件是齐次的、边界是周期的还是退化的，就决定了谱的样子。
非齐次的数据只用来定系数，不参与定谱——这就是 \S6.7 第二条 Remark 的确切含义。

\subhdr{8.4\quad 与前三节的衔接}

\begin{center}
\tabtop
\begin{tabular}{@{}L{0.12\linewidth}L{0.24\linewidth}L{0.52\linewidth}@{}}
\thead{本节} & \thead{对应前面} & \thead{差别} \\
\tabmidrule
\S1 球坐标算子 & §8.2 极坐标算子 & 多一个 $\theta$ 项；$r$ 的公共因子仍是 $1/r^2$ \\
\S3 周期条件 & §8.2 定理 1 & 同一个论证，变量换成 $\varphi$ \\
\S4 Legendre & §8.2 欧拉方程 & 从「常系数/欧拉」升级成「变系数」，靠有界性定谱 \\
\S5 Euler 径向 & §8.2 的 $r^{\pm m}$ & 指数从 $\pm m$ 变成 $l,-(l+1)$ \\
\S7 球贝塞尔 & §8.3 柱贝塞尔 & 方程多一个 $2/r$，换元两步，阶数变成半整数 \\
\end{tabular}
\tabbot
\end{center}

%=====================================================================
\hdr{\S9}{速查与易错}
\phantomsection\label{sec:s9}

\subhdr{9.1\quad 六个式子}

\begin{enumerate}[label=\textbf{(\arabic*)}]
  \item 球坐标 Laplace 算子：
        $\nabla^2u=\dfrac{1}{r^2}\partial_r(r^2\partial_ru)+\dfrac{1}{r^2\sin\theta}\partial_\theta(\sin\theta\,\partial_\theta u)+\dfrac{1}{r^2\sin^2\theta}\partial_\varphi^2u$。
  \item 两次分离后：$\lambda=l(l+1)$，$\mu=m^2$；径向 $r^2R''+2rR'-l(l+1)R=0$。
  \item 关联 Legendre 方程：$\dfrac{d}{dx}\Big[(1-x^2)y'\Big]+\Big(l(l+1)-\dfrac{m^2}{1-x^2}\Big)y=0$，$x=\cos\theta$。
  \item $P_l$ 前几个：$P_0=1$，$P_1=x$，$P_2=(3x^2-1)/2$；模方 $\int_{-1}^1P_l^2dx=2/(2l+1)$。
  \item 球坐标 Laplace 通解：$\sum_{l,m}(A_l^mr^l+C_l^mr^{-(l+1)})P_l^m(\cos\theta)(\cos m\varphi,\sin m\varphi)$。
  \item 球贝塞尔：$R=Cj_l(kr)+Dn_l(kr)$，$j_0=\sin x/x$，$n_0=-\cos x/x$；$k$ 由 $j_l(kr_0)=0$ 定。
\end{enumerate}

\subhdr{9.2\quad 六条易错（每条回指某一步）}

\begin{enumerate}[label=\textbf{\arabic*.}]
  \item \textbf{把 $\theta$ 当成与 $xy$ 平面的夹角。}$\theta$ 自 $+z$ 轴起算，$x=\cos\theta=z/r$（\S1.1）。写错这一条，
        后面 $P_1=z/r$ 全反号。
  \item \textbf{换元时丢负号或写错 $\sin^2\theta$。}$d/d\theta=-\sin\theta\,d/dx$，$\sin^2\theta=1-x^2$（\S4.1 的警示框）。
  \item \textbf{把 $P_l^m$ 的 $m$ 取到无穷。}$m\le l$，$m>l$ 时 $P_l^m\equiv0$（\S4.5 第 (1) 条）。
  \item \textbf{球内问题保留 $r^{-(l+1)}$。}$r=0$ 处发散，必须丢（\S5.2）；球外若远处有外场，$r^l$ 又不能丢。
  \item \textbf{正交性漏掉 $\sin\theta$ 或 $r^2$。}$\theta$ 的积分要乘 $\sin\theta$，$r$ 的积分要乘 $r^2$（\S4.4、\S7.5）。
  \item \textbf{把 $j_l$ 的零点当成 $n\pi$。}只有 $l=0$ 是 $n\pi$；$l\ge1$ 的零点是 $\tan x=x$ 一类的超越方程根（\S7.4）。
\end{enumerate}

\subhdr{9.3\quad 一句话对照 §8.2}

极坐标的四步是「分离 $\to$ 定谱 $\to$ 叠加 $\to$ 定系数」；球坐标多一步「第二次分离」，
其余四步不变，只是特征函数从 $\cos m\varphi$ 变成 $P_l^m(\cos\theta)\cos m\varphi$。
换句话说：\textbf{球坐标不是新方法，是同一个方法在二维球面上的版本}。

%=====================================================================
\hdr{\S10}{三句话总结}
\phantomsection\label{sec:s10}

\begin{enumerate}[label=\textbf{(\arabic*)}]
  \item \textbf{球坐标的两次分离各出一个常数}：第一次（$r$ 与角向）出 $\lambda$，第二次（$\varphi$ 与 $\theta$）出 $\mu=m^2$。
        $\mu$ 由周期条件定，$\lambda$ 由 $x=\pm1$ 的有界性定成 $l(l+1)$。
  \item \textbf{三类条件各管一件事}：周期条件定 $m$，自然边界定 $l$，球面齐次条件定 $k$。
        三者都只定谱，非齐次数据只定系数。
  \item \textbf{球坐标的解 = 球谐 × 径向因子}：Laplace 的径向是 $r^l$ 与 $r^{-(l+1)}$，
        波与输运的径向是 $j_l(kr)$ 与 $n_l(kr)$；球内丢掉在 $r=0$ 发散的那一支，球外丢掉在 $r\to\infty$ 发散的那一支。
\end{enumerate}

\clearpage

% ===================== 练习 =====================
\phantomsection\label{sec:practice}
\hdr{\S11}{独立练习（五题）}

五题分别练：最低阶的定系数（P1）、Legendre 展开（P2）、球壳的两个径向条件（P3）、球对称的球贝塞尔（P4）、
带 $\varphi$ 依赖的球谐（P5）。题干只给条件与考点，步骤全在附录，题号对应。

\practice{P1}{自编，仿 §6.6 例 1}{半径 $a$ 的球内解 $\nabla^2u=0$，球面 $u(a,\theta)=U_0P_2(\cos\theta)$。求 $u(r,\theta)$，并指出球内 $|u|$ 的最大值出现在哪里。}{$l=2$ 的径向因子是 $r^2$；球内只留 $r^l$}{ans:P1}

\practice{P2}{自编，Legendre 展开}{把 $f(\theta)=\cos^3\theta$ 按 $P_l(\cos\theta)$ 展开，即求系数 $c_l$ 使 $f=\sum c_lP_l$。(1) 用 $P_1,P_3$ 的显式直接分解；(2) 用正交性公式验算 $c_1,c_3$，并说明为什么 $c_0=c_2=c_4=\cdots=0$。}{$x^3$ 用 $P_1,P_3$ 表示；奇偶性筛掉偶 $l$}{ans:P2}

\practice{P3}{自编，球壳}{球壳 $a<r<b$ 内解 $\nabla^2u=0$，内球面 $u(a,\theta)=0$，外球面 $u(b,\theta)=U_0\cos\theta$。求 $u(r,\theta)$，并核对两个球面条件。}{球壳两组径向解都留；$l=1$ 的两个系数由两个球面条件定}{ans:P3}

\practice{P4}{自编，球对称输运}{半径 $R$ 的球内，初始温度 $u(r,0)=U_0$（常数），球面保持零温 $u(R,t)=0$。求球对称解 $u(r,t)$ 的级数，并写出前两个非零项。}{$l=0$ 的球贝塞尔是 $j_0=\sin x/x$；$k_n=n\pi/R$；$\int_0^R j_0^2r^2dr$ 的模方}{ans:P4}

\practice{P5}{自编，带 $\varphi$ 依赖}{半径 $a$ 的球内解 $\nabla^2u=0$，球面 $u(a,\theta,\varphi)=U_0\sin\theta\cos\varphi$。求 $u$，并把它写成直角坐标形式。}{$P_1^1(\cos\theta)=\sin\theta$；$r\sin\theta\cos\varphi=x$}{ans:P5}

\clearpage

% ===================== 附录：练习详解 =====================
\phantomsection\label{sec:app}
\hdr{附录 A}{五道练习的详解}

每块先复述题干与来源，再展开。做完一题再翻一题。

\answerstart{P1}{自编，$l=2$ 的球内问题}{半径 $a$ 的球内 $\nabla^2u=0$，$u(a,\theta)=U_0P_2(\cos\theta)$。求 $u$，并指出 $|u|$ 的最大值位置。}

\textbf{第 1 步（认阶）。}球面条件只含 $P_2(\cos\theta)$，即 $l=2$，且与 $\varphi$ 无关（$m=0$）。
所以解只含 $l=2$、$m=0$ 那一项。

\textbf{第 2 步（径向因子）。}球内丢掉 $r^{-(l+1)}$，$l=2$ 时只剩 $r^2$。设
\[
u=C\,r^2P_2(\cos\theta).
\]

\textbf{第 3 步（定系数）。}在 $r=a$ 上：
\[
C\,a^2P_2(\cos\theta)=U_0P_2(\cos\theta)
\quad\Longrightarrow\quad
C=\frac{U_0}{a^2}.
\]
所以
\[
u(r,\theta)=U_0\Big(\frac{r}{a}\Big)^2P_2(\cos\theta)
=\frac{U_0}{2}\Big(\frac{r}{a}\Big)^2(3\cos^2\theta-1).
\]

\textbf{第 4 步（最大值位置）。}$r\le a$，所以 $(r/a)^2\le1$；$P_2(x)=\frac12(3x^2-1)$ 在 $x=\pm1$ 取最大值 $1$，
在 $x=0$ 取最小值 $-\frac12$。两者相乘：
\[
|u|\le U_0,\qquad
|u|=U_0\ \text{当且仅当}\ r=a,\ \theta=0\ \text{或}\ \pi .
\]
即最大值出现在球面的两个极点（$+z$ 与 $-z$ 方向）。

\textbf{核对。}①$\nabla^2u=0$：$r^2P_2(\cos\theta)$ 是 $z$ 的二次调和多项式（等于 $z^2-r^2/3$ 的常数倍），调和 $\checkmark$；
②球面取值：$r=a$ 时回到 $U_0P_2$ $\checkmark$；③$r=0$ 有界 $\checkmark$；④赤道 $\theta=\pi/2$ 上 $u=-\frac{U_0}{2}(r/a)^2<0$，与 $P_2$ 的符号一致 $\checkmark$。

\textbf{直角坐标形式}（可选）：$P_2=(3z^2-r^2)/(2r^2)$，所以
\[
u=\frac{U_0}{2a^2}(3z^2-r^2).
\]

\answerstart{P2}{自编，Legendre 展开}{把 $f(\theta)=\cos^3\theta$ 按 $P_l$ 展开，求系数并验算。}

\textbf{第 1 步（直接分解）。}令 $x=\cos\theta$，要分解 $x^3$。由
\[
P_1=x,\qquad P_3=\frac{5x^3-3x}{2}
\]
反解 $x^3$：
\[
x^3=\frac{2P_3+3P_1}{5}=\frac{3}{5}P_1+\frac{2}{5}P_3 .
\]
所以 $c_1=\frac35$、$c_3=\frac25$，其余系数为零。

\textbf{第 2 步（正交性验算）。}公式
\[
c_l=\frac{2l+1}{2}\int_{-1}^{1}x^3P_l(x)\,dx .
\]
对 $l=1$：
\[
c_1=\frac{3}{2}\int_{-1}^1x^3\cdot x\,dx=\frac32\cdot\frac{2}{5}=\frac35 .
\]
对 $l=3$：
\[
c_3=\frac{7}{2}\int_{-1}^1x^3\cdot\frac{5x^3-3x}{2}\,dx
=\frac74\Big(5\cdot\frac{2}{7}-3\cdot\frac{2}{5}\Big)
=\frac74\Big(\frac{10}{7}-\frac65\Big)
=\frac74\cdot\frac{8}{35}=\frac25 .
\]
两式与第 1 步一致。

\textbf{第 3 步（为什么偶 $l$ 为零）。}$x^3$ 是奇函数；$P_l$ 的奇偶性与 $l$ 相同（$l$ 偶则偶）。
奇函数乘偶函数是奇函数，在 $[-1,1]$ 上积分为零，所以 $c_0=c_2=c_4=\cdots=0$。
$l>3$ 时也不需要算：$x^3$ 是三次多项式，与次数更高的 $P_l$ 正交（或直接用正交性公式，积分的被积函数次数不够）。

\textbf{一句话}：\textbf{对称性先筛掉一半系数，显式分解再定剩下的}——这是 Legendre 展开最快的做法。

\answerstart{P3}{自编，球壳 $a<r<b$}{内球面 $u(a)=0$，外球面 $u(b)=U_0\cos\theta$，求 $u$。}

\textbf{第 1 步（认阶）。}边界含 $P_1(\cos\theta)$，故 $l=1$、$m=0$。

\textbf{第 2 步（径向因子）。}球壳上 $r^1$ 与 $r^{-2}$ 都有界，两组都留：
\[
u(r,\theta)=\Big(Cr+\frac{D}{r^2}\Big)\cos\theta .
\]

\textbf{第 3 步（两个球面条件）。}
\[
\begin{cases}
Ca+\dfrac{D}{a^2}=0,\\[6pt]
Cb+\dfrac{D}{b^2}=U_0 .
\end{cases}
\]
由第一式 $D=-Ca^3$，代入第二式：
\[
C\Big(b-\frac{a^3}{b^2}\Big)=U_0
\quad\Longrightarrow\quad
C=\frac{U_0b^2}{b^3-a^3},
\]
于是
\[
D=-\frac{U_0a^3b^2}{b^3-a^3}.
\]

\textbf{第 4 步（答案）。}
\[
u(r,\theta)=\frac{U_0b^2}{b^3-a^3}\Big(r-\frac{a^3}{r^2}\Big)\cos\theta .
\]

\textbf{核对。}①$r=a$：$a-\frac{a^3}{a^2}=0$ $\checkmark$；②$r=b$：$\frac{U_0b^2}{b^3-a^3}\cdot\frac{b^3-a^3}{b^2}=U_0$ $\checkmark$；
③$a\to0$（退化成实心球）：$C\to U_0/b$、$D\to0$，$u\to U_0\frac{r}{b}\cos\theta$，与 §6.6 例 1 一致 $\checkmark$；
④$\nabla^2u=0$：两项分别是 $z$ 与 $z/r^3$ 的常数倍，都调和 $\checkmark$。

\answerstart{P4}{自编，球对称输运}{半径 $R$ 的球内，$u(r,0)=U_0$，$u(R,t)=0$。求球对称解。}

\textbf{第 1 步（球对称的方程）。}$u=u(r,t)$，Laplace 算子只剩径向：
\[
u_t=a^2\frac{1}{r^2}\frac{\partial}{\partial r}\Big(r^2\frac{\partial u}{\partial r}\Big).
\]

\textbf{第 2 步（分离）。}设 $u=R(r)T(t)$：
\[
\frac{T'}{a^2T}=\frac{(r^2R')'}{r^2R}=-k^2
\quad\Longrightarrow\quad
T'+a^2k^2T=0,\qquad
(r^2R')'+k^2r^2R=0 .
\]
这是 $l=0$ 的球贝塞尔方程。

\textbf{第 3 步（径向解与 $k$）。}$l=0$ 时 $R=Cj_0(kr)+Dn_0(kr)$，$n_0$ 在 $r=0$ 发散，丢掉。$j_0(x)=\sin x/x$，所以
\[
R=C\,\frac{\sin(kr)}{kr}.
\]
球面零温：$\sin(kR)=0$，得
\[
k_n=\frac{n\pi}{R},\qquad n=1,2,3,\dots
\]

\textbf{第 4 步（时间因子与叠加）。}$T=e^{-a^2k^2t}$，故
\[
u(r,t)=\sum_{n=1}^{\infty}A_n\,\frac{\sin(k_nr)}{k_nr}\,e^{-a^2k_n^2t}.
\]

\textbf{第 5 步（定系数）。}用球贝塞尔的正交性（权 $r^2$）：
\[
\int_0^Rj_0(k_nr)j_0(k_mr)\,r^2dr=\frac{R^3}{2n^2\pi^2}\delta_{nm}
\]
（一般公式是 $\frac{R^3}{2}j_{l+1}(k_nR)^2$，$l=0$、$k_nR=n\pi$ 时 $j_1(n\pi)=(-1)^{n+1}/(n\pi)$，平方即 $1/(n^2\pi^2)$）。
于是
\[
A_n=\frac{2n^2\pi^2}{R^3}\int_0^RU_0\,\frac{\sin(k_nr)}{k_nr}r^2dr
=\frac{2n^2\pi^2U_0}{R^3k_n}\int_0^Rr\sin(k_nr)\,dr .
\]
分部积分：
\[
\int_0^Rr\sin(k_nr)\,dr=\Big[-\frac{r\cos(k_nr)}{k_n}+\frac{\sin(k_nr)}{k_n^2}\Big]_0^R
=-\frac{R(-1)^n}{k_n}.
\]
代入并用 $k_n=n\pi/R$：
\[
A_n=\frac{2n^2\pi^2U_0}{R^3k_n}\cdot\Big(-\frac{R(-1)^n}{k_n}\Big)
=\frac{2U_0(-1)^{n+1}}{1}=2U_0(-1)^{n+1}.
\]

\textbf{第 6 步（答案与前两项）。}
\[
u(r,t)=2U_0\sum_{n=1}^{\infty}(-1)^{n+1}\frac{\sin(n\pi r/R)}{n\pi r/R}\,e^{-a^2n^2\pi^2t/R^2}.
\]
前两项：
\[
u=2U_0\frac{\sin(\pi r/R)}{\pi r/R}e^{-a^2\pi^2t/R^2}
-2U_0\frac{\sin(2\pi r/R)}{2\pi r/R}e^{-4a^2\pi^2t/R^2}+\cdots .
\]

\textbf{核对。}①$r=R$：$\sin(n\pi)=0$，$u(R,t)=0$ $\checkmark$；②$t=0$：用 $\sum_{n\ge1}(-1)^{n+1}\frac{\sin nx}{n}=\frac{x}{2}$（$0<x<\pi$），
令 $x=\pi r/R$ 得 $u(r,0)=2U_0\cdot\frac{\pi r/R}{2}\big/\frac{\pi r}{R}=U_0$ $\checkmark$；
③$t\to\infty$：首项指数最慢，$u\to0$，球面零温、热量不断流出，合理 $\checkmark$；
④$r=0$ 有界：$\sin(k_nr)/(k_nr)\to1$ $\checkmark$。

\answerstart{P5}{自编，带 $\varphi$ 依赖}{半径 $a$ 的球内 $\nabla^2u=0$，$u(a,\theta,\varphi)=U_0\sin\theta\cos\varphi$。求 $u$ 与直角坐标形式。}

\textbf{第 1 步（认阶）。}边界里 $\sin\theta\cos\varphi$ 可拆成
\[
\sin\theta\cos\varphi=P_1^1(\cos\theta)\cos\varphi
\]
（用了 $P_1^1(x)=(1-x^2)^{1/2}=\sin\theta$，见 \eqref{eq:P11}）。所以 $l=1$、$m=1$，只含 $\cos\varphi$ 那一支。

\textbf{第 2 步（设解）。}球内只留 $r^l=r$：
\[
u=C\,rP_1^1(\cos\theta)\cos\varphi=C\,r\sin\theta\cos\varphi .
\]

\textbf{第 3 步（定系数）。}在 $r=a$ 上：
\[
C\,a\sin\theta\cos\varphi=U_0\sin\theta\cos\varphi
\quad\Longrightarrow\quad
C=\frac{U_0}{a}.
\]
所以
\[
u(r,\theta,\varphi)=\frac{U_0}{a}\,r\sin\theta\cos\varphi .
\]

\textbf{第 4 步（直角坐标形式）。}因为 $x=r\sin\theta\cos\varphi$，
\[
u=\frac{U_0}{a}x .
\]
这是沿 $x$ 方向的均匀场。

\textbf{核对。}①$\nabla^2x=0$ $\checkmark$；②球面 $r=a$：$u=U_0\sin\theta\cos\varphi$ $\checkmark$；③$r=0$ 有界 $\checkmark$；
④与 P1 对照：P1 的边界是 $P_2$（$z^2$ 型），P5 的是 $P_1^1\cos\varphi$（$x$ 型）——两题的解分别是 $z^2-r^2/3$ 型与 $x$ 型，
次数与直角坐标的对应关系一致（$l$ 次 → $l$ 次多项式）。

\flowerrule

\textbf{五题合起来看}：P1、P3、P5 练定系数（球内、球壳、带 $\varphi$），P2 练展开，P4 练球贝塞尔。
四件事的共同骨架仍是 \S0 的六步：分离 $\to$ 定谱（$m$、$l$、$k$）$\to$ 按区域取舍 $\to$ 叠加 $\to$ 正交定系数 $\to$ 核对。
把 P4 的 $j_l(kr)$ 换成 $r^l$、把 $e^{-a^2k^2t}$ 换成 $1$，它就退化成 P1 的结构——这正是「同一个方法」的意思。




\end{document}


